%Document Class
\documentclass[a4paper,12pt]{article}


%Packages 
\usepackage{amscd,amssymb,amsmath,amsfonts,graphicx,verbatim,hyperref}
\usepackage[OT1,T1]{fontenc}
\usepackage[all]{xypic}
\usepackage{pdfpages}


%Theoretic Environment Setup
\newtheorem{theorem}{Theorem}[section]
\newtheorem{corollary}[theorem]{Corollary}
\newtheorem{lemma}[theorem]{Lemma}
\newtheorem{proposition}[theorem]{Proposition}
\newtheorem{remark}[theorem]{Remark}
\newtheorem{definition}[theorem]{Definition}
\newtheorem{example}[theorem]{Example}
\newenvironment{proof}[1][Proof]{\begin{trivlist}
\item[\hskip \labelsep {\bfseries #1}]}{\end{trivlist}}


%Tom's Not-So-Ridiculous Shortcuts
\newcommand{\oop}{\operatorname{op}}
\newcommand{\op}{^{\oop}}
\newcommand{\pps}{\operatorname{ps}}
\newcommand{\ps}{^{\pps}}
\newcommand{\lowps}{_{\pps}}
\newcommand{\ttract}{\operatorname{tract}}
\newcommand{\tract}{_{\ttract}}
\newcommand{\rrep}{\operatorname{rep}}
\newcommand{\rep}{_{\rrep}}
\newcommand{\bicat}{\mathrm{Bicat}(\bc\op,\bm\op)\op}
\newcommand{\lom}{\ell{om}}
\newcommand{\rom}{r{om}}
\newcommand{\warmon}{\,\Box\,}
\numberwithin{equation}{section}


%French Greek Shortcuts
\def\th{\theta}
\def\Th{\Theta}
\def\ga{\gamma}
\def\Ga{\Gamma}
\def\be{\beta}
\def\al{\alpha}
\def\ep{\epsilon}
\def\la{\lambda}
\def\La{\Lambda}
\def\de{\delta}
\def\De{\Delta}
\def\om{\omega}
\def\Om{\Omega}
\def\sig{\sigma}
\def\Sig{\Sigma}
\def\vphi{\varphi}


%French Calligraphic Character Shortcuts 
\def\CA{{\cal A}}
\def\CB{{\cal B}}
\def\CC{{\cal C}}
\def\CD{{\cal D}}
\def\CE{{\cal E}}
\def\CF{{\cal F}}
\def\CG{{\cal G}}
\def\CH{{\cal H}}
\def\CI{{\cal J}}
\def\CJ{{\cal J}}
\def\CK{{\cal K}}
\def\CL{{\cal L}}
\def\CM{{\cal M}}
\def\CN{{\cal N}}
\def\CO{{\cal O}}
\def\CP{{\cal P}}
\def\CQ{{\cal Q}}
\def\CR{{\cal R}}
\def\CS{{\cal S}}
\def\CT{{\cal T}}
\def\CU{{\cal U}}
\def\CV{{\cal V}}
\def\CW{{\cal W}}
\def\CX{{\cal X}}
\def\CY{{\cal Y}}
\def\CZ{{\cal Z}}

%Ross' Shortcuts 

\def\lra{{\longrightarrow}}
\def\Lra{{\Longrightarrow}}
\def\lla{{\longleftarrow}}
\def\Lla{{\Longleftarrow}}


%French LaTeX Oddities
\def\qq{ \begin{eqnarray} }
\def\qqq{ \end{eqnarray} }
\def\non{ \nonumber }

\newcommand{\no}{\noindent}
\newcommand{\vs}{\vspace}
\newcommand{\hs}{\hspace}
\newcommand{\e}{�}
\newcommand{\ef}{�}
\newcommand{\p}{\partial}
\newcommand{\un}{\underline}
\newcommand{\ha}{{1\over 2}}
\newcommand{\pa}{\partial}
\newcommand{\sh}{\sinh}
\newcommand{\ch}{\cosh}
\newcommand{\rea}{{\rm Re}}
\newcommand{\ima}{{\rm Im}}

\def\theequation{{\thesection.\arabic{equation}}}


%French Document Margins
\setlength{\textwidth}{16cm}
\setlength{\textheight}{23.315cm}
\hoffset -1.2cm
\topmargin= -1cm
\raggedbottom
\renewcommand{\baselinestretch}{1.5}
\pagestyle{plain}


%Begin Document Body

\begin{document}

\author{Ross Street\footnote{The author gratefully acknowledges the support of an Australian Research Council Discovery Grant DP1094883.}}
\title{Categories in categories, and size matters}
\date{}
\maketitle

\noindent {\small{\emph{2010 Mathematics Subject Classification} \quad 18D10; 18D20; 18D35}}
\\
{\small{\emph{Key words and phrases.} internal category; category object; internal full subcategory; fibration; double category; descent category; distributors; profunctors; modules.}}

\begin{abstract}
\noindent We look again at internal categories and internal full subcategories in a category $\CC$. We look at the relationship between a generalised Yoneda lemma and the descent construction. Application to $\CC = \mathrm{Cat}$ gives results on double categories.

\end{abstract}



\section*{Introduction}

We revisit the theory of categories in a category.  We look in particular at categories in the category of categories; that is, at double categories.

\tableofcontents

\section{Categories in categories}\label{Cic}

A \textit{category} $C$ \textit{in} a category $\CC$ is a diagram 
\begin{equation}
\xymatrix{
C : & C_{3} \ar@<+3ex>[rr]|-{d_{0}} \ar@<+1ex>[rr]|-{d_{1}} \ar@<-1ex>[rr]|-{d_{2}} \ar@<-3ex>[rr]|-{d_{3}} && \ar@<+4ex>[rr]|-{d_{0}} \ar@{<-}@<+2ex>[rr]|-{i_{0}} \ar[rr]|-{d_{1}} \ar@{<-}@<-2ex>[rr]|-{i_{1}} \ar@<-4ex>[rr]|-{d_{2}} C_{2} && C_{1} \ar@<+2ex>[rr]|-{d_{0}} \ar@{<-}[rr]|-{i} \ar@<-2ex>[rr]|-{d_{1}} && C_{0}
}
\end{equation}
in $\CC$ such that
\begin{itemize}
\item[C1.] $d_p d_q = d_{q-1} d_p$ for $p<q$
\item[C2.] $d_0 i = d_1 i = 1_{C_0},   d_0 i_0 = d_1 i_0 = d_1 i_1 = d_2 i_1 = 1_{C_1}$
\item[C3.] $i_0 i= i_1 i, d_2 i_0 = i d_1, d_0 i_1 = i d_0$
\item[C4.] the following two squares are pullbacks.
\begin{equation}\label{pbc}
\xymatrix{
C_2 \ar[rr]^-{d_0} \ar[d]_-{d_2} && C_1 \ar[d]^-{d_1} \\
C_1 \ar[rr]_-{d_0} && C_0}\qquad
\xymatrix{
C_3 \ar[rr]^-{d_0} \ar[d]_-{d_3} && C_2 \ar[d]^-{d_2} \\
C_2 \ar[rr]_-{d_0} && C_1}
\end{equation}
\end{itemize}
The definition goes back to Ehresmann \cite{Ehres1963}.

For any object $U$ of $\CC$ and any category $C$ in $\CC$, we obtain a category $\CC (U,C)$ as follows. An object is a morphism $u:U\lra C_0$ in $\CC$. A morphism $\gamma : u \lra v$ is a morphism $\gamma : U \lra C_1$ in $\CC$ such that $d_0 \gamma = u$ and $d_1 \gamma = v$. The identity $1_u : u \lra u$ is $iu : U \lra C_1$. The composite of $\gamma : u \lra v$ and $\delta : v \lra w$ is $d_1 (\ga , \de) : U \lra C_1$ where $(\ga , \de):U\lra C_2$ is defined by $d_0 (\ga , \de) = \ga$ and $d_2 (\ga , \de) =\de$.

For each morphism $h:V\lra U$ in $\CC$, we have a functor $$\CC (h,C): \CC (U,C) \lra \CC (V,C)$$ taking $u$ to $uh$ and $\ga$ to $\ga h$. Thus we obtain a functor $$\CC (-,C): \CC^{\mathrm{op}} \lra \mathrm{Cat} \ .$$

A \textit{functor} $f:C\lra D$ between categories $C$ and $D$ in $\CC$ is a morphism of diagrams 
$$\xymatrix{
\ar@<+2ex>[rr]|-{d_{0}} \ar[rr]|-{d_{1}} \ar@<-2ex>[rr]_(0){\phantom{.}}="2a"|-{d_{2}} C_{2} \ar@{}[ddrr]|-{ } && C_{1} \ar@<+2ex>[rr]|-{d_{0}} \ar@{<-}[rr]|-{i} \ar@<-2ex>[rr]_(0){\phantom{.}}="1a"_(1){\phantom{.}}="0a"|-{d_{1}} \ar@{}[ddrr]|-{ } && C_{0}  \\ \\
\ar@<+2ex>[rr]^(0){\phantom{.}}="2b"|-{d_{0}} \ar[rr]|-{d_{1}} \ar@<-2ex>[rr]|-{d_{2}} D_{2} && D_{1} \ar@<+2ex>[rr]^(0){\phantom{.}}="1b"^(1){\phantom{.}}="0b"|-{d_{0}} \ar@{<-}[rr]|-{i} \ar@<-2ex>[rr]|-{d_{1}} && D_{0} 
\ar"2a";"2b"_{f_{2}} \ar"1a";"1b"_{f_{1}} \ar"0a";"0b"_{f_{0}}
}$$

Write $\mathrm{Cat} \CC$ for the category of categories in $\CC$ and functors between them.

Each functor $f:C\lra D$ defines a functor $$\CC (U,f):\CC (U,C) \lra \CC (U,D)$$ taking $u$ to $f_0 u$ and $\ga$ to $f_1 \ga$. Thus we obtain a functor $$\CC (-,-) : \CC^{\mathrm{op}} \times \mathrm{Cat} \CC \lra \mathrm{Cat}$$ which can be recast as a functor
\begin{equation}\label{eYe}
\mathrm{Yon} : \mathrm{Cat} \CC \lra [\CC^{\mathrm{op}} , \mathrm{Cat}]
\end{equation}
extending the usual Yoneda embedding $\CC \lra [\CC^{\mathrm{op}} , \mathrm{Set}] $.

A \textit{natural transformation} $\theta : f \Lra g : C \lra D$ between functors $f$ and $g$ is a morphism $$\theta : C_0 \lra D_1$$ in $\CC$ such that
\begin{itemize}
\item[N1.] $d_0 \theta = f_0$ and $d_1 \theta = g_0$
\item[N2.] the following diagram commutes.
\begin{equation}
\xymatrix{
 & C_{1}  \ar@<+1ex>[rr]|-{(f_{1}, \theta d_{1})} \ar@<-1ex>[rr]|-{(\theta d_{0},g_{1})}  &&  \ar[rr]|-{d_{1}}  D_{2} && D_{1} 
}
\end{equation}
\end{itemize}

Given a diagram

\begin{equation}
 \xymatrix{B \ar[rr]^r  && C  \ar@/^1pc/[rr]^f 
\ar@/_1pc/[rr]_{g} \ar@{}[rr]|{\Downarrow \; \theta}&& D  \ar[rr]^s && E}
\end{equation}
of functors $r$, $f$, $g$, $s$ and natural transformation $\theta$, we obtain a `whiskered' natural transformation $s\theta r : sfr \Longrightarrow sgr$ defined as the composite
\begin{equation}
 \xymatrix{B_0 \ar[rr]^{r_{0}}  && C_0 \ar[rr]^{\theta} && D_1  \ar[rr]^{s_{1}} && E_1 \ \ .}
\end{equation}
Given a diagram

\begin{equation}
 \xymatrix{\ar[rr]|{g} C \ar@/^1.5pc/[rr]^{f}_{\Downarrow \; \theta}  \ar@/_1.5pc/[rr]_{h}^{\Downarrow \; \phi} & & D 
 }
\end{equation}
of functors $f$, $g$, $h$ and natural transformations $\theta$, $\phi$, we obtain a natural transformation $\phi \circ \theta : f \Longrightarrow h$ defined as the composite
\begin{equation}
 \xymatrix{C_0 \ar[rr]^{(\theta , \phi)}  && D_2 \ar[rr]^{d_1} && D_1  \ \ .}
\end{equation}
In this way, $\mathrm{Cat} \CC$ becomes a 2-category.

Each object $U$ of $\CC$ and natural transformation $\theta:f \Lra g:C \lra D$ define a natural transformation
 $$
 \xymatrix{\CC(U,C) \ar@/^1pc/[rr]^{\CC(U,f)}  \ar@/_1pc/[rr]_{\CC(U,g)}  \ar@{}[rr]|{\Downarrow \; \CC(U,\theta)}&& \CC(U,D) 
 }$$
 whose component at $u \in \CC(U,C)$ is $\theta u : fu\lra gu$ in $\CC(U,D)$. In this way, we see that the extended Yoneda embedding \eqref{eYe} is a 2-functor.
 
For each category $C$ in $\CC$, there is a category $\int C$ defined as follows. The objects are pairs $(U,u)$ where $U$ is an object of $\CC$ and $u$ is an object of $\CC(U,C)$. A morphism $(h,\theta):(U,u)\lra (V,v)$ consists of a morphism $h:U\lra V$ in $\CC$ and a morphism $\theta :u\lra vh$ in $\CC(U,C)$. Regarding each object $U$ of $\CC$ as a `discrete category' in $\CC$ (that is, all $U_n =U$ and $d_p =1_U$), we can say that morphisms of $\int C$ are diagrams

\begin{equation}
\xymatrix{
U \ar[rd]_{u}^(0.5){\phantom{a}}="1" \ar[rr]^{h}  && V \ar[ld]^{v}_(0.5){\phantom{a}}="2" \ar@{=>}"1";"2"^-{\theta}
\\
& C 
}
\end{equation}
in $\mathrm{Cat}\CC$; composition in $\int C$ is by pasting these triangles. There is a functor
\begin{equation}\label{Pfromint} 
P : \begingroup\textstyle\int\endgroup C \lra \CC
\end{equation}
defined by $P(U,u)=U$ and $P(h,\theta)=h$.

\section{Fibrations and the generalized Yoneda Lemma}\label{FgYL}

Consider a functor $P:\CX\lra\CC$. A morphism $x:Y\lra X$ in $\CX$ is \textit{cartesian} (with respect to $P$) when the following square is a pullback for all objects $K$ of $\CX$.
\begin{equation}
\xymatrix{
\CX(K,Y) \ar[rr]^-{\CX(K,x)} \ar[d]_-{P} && \CX(K,X) \ar[d]^-{P} \\
\CX(PK,PY) \ar[rr]_-{\CX(PK,Px)} && \CX(PK,PX)}
\end{equation}

\begin{lemma}\label{cart&pb} Consider a commutative square
$$
\xymatrix{
Y^{\prime} \ar[rr]^-{x^{\prime}} \ar[d]_-{k} && X^{\prime} \ar[d]^-{h} \\
Y \ar[rr]_-{x} && X}
$$
in $\CX$, which is taken to a pullback in $\CC$ by the functor $P:\CX\lra\CC$, and where $x$ is cartesian. The square is a pullback if and only if $x^{\prime}$ is cartesian. 

\end{lemma}

\begin{proof} Contemplate the two diagrams below.
$$
\xymatrix{
\CX(K,Y^{\prime}) \ar[rr]^-{\CX(1,x^{\prime})} \ar[d]_-{\CX(1,k)} && \CX(K,X^{\prime}) \ar[d]^-{\CX(1,h)} \\
\CX(K,Y)\ar[rr]_-{\CX(1,x)} \ar[d]_-{P} && \CX(K,X) \ar[d]^-{P} \\
\CC(PK,PY) \ar[rr]_-{\CC(1,Px)} && \CC(PK,PX)}\qquad
\xymatrix{
\CX(K,Y^{\prime}) \ar[rr]^-{\CX(1,x^{\prime})} \ar[d]_-{P} && \CX(K,X^{\prime}) \ar[d]^-{P} \\
\CC(PK,PY^{\prime}) \ar[rr]_-{\CC(1,Px^{\prime})} \ar[d]_-{\CC(1,Pk)} && \CC(PK,PX^{\prime}) \ar[d]^-{\CC(1,Ph)} \\
\CC(PK,PY) \ar[rr]_-{\CC(1,Px)} && \CC(PK,PX}
$$
The two squares obtained by composing both diagrams vertically are equal. By hypothesis, the bottom square in each diagram is a pullback. It follows that the top square in one diagram is a pullback if and only if the top square in the other is.  \hfill{$\square$}

\end{proof}

For any functor $P:\CX\lra\CC$ and any category $C$ in $\CC$, we define a category $\CX^C$. The objects are categories $X$ in $\CX$ such that $PX=C$ and $d_0 :X_1 \lra X_0$ is cartesian. A morphism $r:X \lra Y$ is a functor in $\CX$ such that $Pr:PX\lra PY$ is the identity functor of $C$.

\begin{remark}\label{pbcart} By Lemma~\ref{cart&pb}, the pullback requirement C4 \eqref{pbc}) on a category $X$ in $\CX$ belonging to $\CX^C$ can be replaced by the requirement that $$d_0:X_2\lra X_1 \ \  \text{and} \ \   d_0:X_3\lra X_2$$ should be cartesian.
\end{remark}

For functors $P:\CX\lra \CC$ and $Q:\CY\lra \CC$, we write 
$$\mathrm{Cat}_{/\CC}(\CX,\CY)$$ 
for the category of functors $F:\CX\lra \CY$ with $QF=P$. The morphisms are natural transformations taken to identities by $Q$. We write
$$\mathrm{Cart}_{/\CC}(\CX,\CY)$$
for the full subcategory of $\mathrm{Cat}_{/\CC}(\CX,\CY)$ consisting of those functors $F$ which preserve cartesian morphisms. Such an $F$ defines a functor $$F^C :\CX^C\lra \CY^C$$ given by $F^CX=FX$; this uses Remark~\ref{pbcart}. We obtain a functor
\begin{equation}
(-)^C: \mathrm{Cart}_{/\CC}(\CX,\CY)\lra[\CX^C,\CY^C]
\end{equation}
for each category $C$ in $\CC$.

A functor $P:\CX\lra \CC$ is called a \textit{fibration} (in the sense of Grothendieck \cite{GrothDesc}, \cite{GrothFond}) when, for all $h:V\lra U$ in $\CC$ and $X$ in $\CX$ with $PX=U$, there exists a cartesian morphism $x:Y\lra X$ with $Px=h$. Such an $x:Y\lra X$ is called an \textit{inverse image} of $X$ along $h$. A choice of inverse image for all such $X$ and $h$ is called a \textit{cleavage} for $P$; we write $$h^X:h^{\star}X\lra X$$ for the chosen cartesian morphisms. 

\noindent {\bf Convention} Given a cleavage, by $\CX^C$ we will then mean the restriction of $\CX^C$ to those objects $X$ with all the morphisms $d_0:X_{n+1}\lra X_n$ taken to be chosen cartesian.  

Let us now return to the functor $P: \int C \lra \CC$ of \eqref{Pfromint}. A morphism $(h,\theta):(U,u)\lra (V,v)$ is cartesian if and only if $\theta :u \lra vh$ is invertible in the category $\CC(U,C)$. In fact, $P: \int C \lra \CC$  is a fibration and there is a cleavage for which the chosen cartesian $(h,\theta)$ have $\theta$ an identity.

The following diagram defines a category $\hat{C}$ in $\int C$:
\begin{equation}
\xymatrix{
& (C_{2}, d_0d_0) \ar@<+2ex>[rr]|-{(d_{0}, id_0d_0)} \ar@<+0ex>[rr]|-{(d_{1},id_0d_0)} \ar@<-2 ex>[rr]|-{(d_{2},d_0)} &&   (C_{1},d_0) \ar@<+2ex>[rr]|-{(d_{0},id_0)} \ar@{<-}[rr]|-{(i,i)} \ar@<-2ex>[rr]|-{(d_{1},1_{C_1})} && (C_{0},1_{C_0}) \ \ . 
}
\end{equation}
For each object $(U,u)$ of $\int C$, we can identify the category
\begin{equation}
\begingroup\textstyle\int\endgroup C((U,u), \hat{C})=u/\CC(U,C)
\end{equation}
as the category of objects of $\CC(U,C)$ under the object $u$.

Note moreover that $\hat C$ is actually an object of $(\int C)^C$. The generalized Yoneda Lemma of \cite{CoIC} Theorem 5.15 expresses the `generic' property of $\hat{C} \in (\int C)^C$.

\begin{theorem}\label{gYl} For any fibration $P:\CX \lra \CC$ and category $C$ in $\CC$, the composite functor 
$$\mathrm{Cart}_{/\CC}(\begingroup\textstyle\int\endgroup C,\CX) \stackrel{(-)^C}{\lra} [(\begingroup\textstyle\int\endgroup C)^C,\CX^C]  \stackrel{\mathrm{eval}_{\hat{C}}}{\lra}\CX^C$$
is an equivalence of categories.
\end{theorem}

\begin{proof} Take $X\in \CX ^C$. We must define a cartesian-morphism-preserving functor
$$\bar{X} : \begingroup\textstyle\int\endgroup C \lra \CX$$
over $\CC$. For $(U,u)\in \begingroup\textstyle\int\endgroup C$, put 
$$\bar{X} (U,u)=u^{\star}X_0 \ \ .$$ 
For $(h,\theta):(U,u)\lra (V,v)$ in $\begingroup\textstyle\int\endgroup C$, notice that the composite of cartesian morphisms
$$\theta^{\star}X_1\stackrel{\theta^{X_1}}{\lra} X_1 \stackrel{d_0}{\lra}X_0$$
is cartesian over $u:U\lra C_0$ (since $d_0 \theta = u$), so we have an isomorphism $u^{\star}X_0 \cong \theta^{\star} X_1$ over $U$ and commuting with $u^{X_0}$ and $d_0\theta^{X_1}$. Define $\bar{X}(h,\theta):\bar{X}(U,u) \lra \bar{X}(V,v)$ such that
$$
\xymatrix{
u^{\star}X_0 \ar[r]^-{\cong} \ar[d]_-{\bar{X}(h,\theta)} & \theta^{\star}X_1 \ar[r]^-{ \theta^{X_1}}  & X_1 \ar[d]^-{d_1} \\
V \ar[rr]_-{v^{X_0}} & & X_0 }
$$
commutes and $P\bar{X}(h,\theta)=h$. Please see \cite{CoIC} Theorem 5.15 for further details.  \hfill{$\square$}

\end{proof}

A \textit{split} fibration is a fibration equipped with a cleavage satisfying the identities
$$ {1_U}^{\star}X=X \ \ , \ \  {1_U}^{X}=1_X \ \ , $$
and
$$(uv)^{\star}X=v^{\star}u^{\star}X \ \ , \ \  (uv)^{X}=u^{X}v^{u^{\star}X} \ \ .$$
Write $\mathrm{Spl}_{/\CC}(\CX,\CY)$ for the full subcategory of $\mathrm{Cart}_{/\CC}(\CX,\CY)$ consisting of those $F:\CX\lra \CY$ which preserve the cleavages given in $\CX$ and $\CY$. Recall our convention about $\CX^C$ in the cloven case. 

Notice that $P: \int C \lra \CC$ is a split fibration. In fact, \cite{CoIC} Theorem 5.15 actually included the following result. Here we make implicit use of our convention.

\begin{proposition}\label{gYlspl} For any split fibration $P:\CX\lra \CC$ and category $C$ in $\CC$, the equivalence of Theorem~\ref{gYl} restricts to an isomorphism of categories $$\mathrm{Spl}_{/\CC}(\begingroup\textstyle\int\endgroup C,\CX) \cong \CX^C \ \ . $$ 
\end{proposition}

\begin{corollary}\label{discretedensity} The 2-functor \eqref{eYe} $$\mathrm{Yon} : \mathrm{Cat} \CC \lra [\CC^{\mathrm{op}} , \mathrm{Cat}]$$ is fully faithful. That is, the inclusion 2-functor $$\CC\lra \mathrm{Cat}\CC$$ is dense.
\end{corollary}

\begin{proof} For categories $C$ and $D$ in $\CC$, we have a commutative square
$$\xymatrix{
(\begingroup\textstyle\int\endgroup D)^C \ar[rr]^-{\cong} \ar[d]_-{\cong} && \mathrm{Spl}_{/\CC}(\begingroup\textstyle\int\endgroup C,\begingroup\textstyle\int\endgroup D) \ar[d]^-{\cong} \\
\mathrm{Cat}\CC(C,D) \ar[rr]_-{\mathrm{Yon}} &&  [\CC^{\mathrm{op}} , \mathrm{Cat}](\CC(-,C),\CC(-,D))}
$$
where the vertical isomorphisms are fairly straightforward and the top isomorphism comes from Proposition~\ref{gYlspl}.  \hfill{$\square$}

\end{proof}

\begin{corollary}\label{discretedetectlim} Limits in the 2-category $\mathrm{Cat}\CC$ are detected by discrete categories in $\CC$.  That is, a (weighted) diagram in $\mathrm{Cat}\CC$ is a limit if and only if the 2-functor \eqref{eYe} takes it to a limit in $[\CC^{\mathrm{op}} , \mathrm{Cat}]$.
\end{corollary}

Perhaps the reader noticed that $d_1:\hat{C}_2\lra \hat{C}_1$ in the generic example is actually cartesian; in fact, it is part of the cleavage. So the following is actually a consequence of Theorem~\ref{gYl}. We give a direct proof.

\begin{proposition} For any object $X$ of $\CX^C$, the morphism $d_1:X_2\lra X_1$ is cartesian.
\end{proposition}

\begin{proof} We have $d_0d_1 = d_0d_0$ and commutative diagrams.
$$
\xymatrix{
\CX(K,X_2) \ar[rr]^-{P} \ar[d]_-{\CX(1,d_1)} && \CC(PK,C_2) \ar[d]^-{\CC(1,d_1)} \\
\CX(K,X_1) \ar[rr]_-{P} \ar[d]_-{\CX(1,d_0)} && \CC(PK,C_1) \ar[d]^-{\CC(1,d_0)} \\
\CX(K,X_0) \ar[rr]_-{P} && \CC(PK,C_0)}\qquad
\xymatrix{
\CX(K,X_2) \ar[rr]^-{P} \ar[d]_-{\CX(1,d_0)} && \CC(PK,C_2) \ar[d]^-{\CC(1,d_0)} \\
\CX(K,X_1) \ar[rr]_-{P} \ar[d]_-{\CX(1,d_0)} && \CC(PK,C_1) \ar[d]^-{\CC(1,d_0)} \\
\CX(K,X_0) \ar[rr]_-{P} && \CC(PK,C_0)}
$$
So the top left square is a pullback, implying $d_1:X_2\lra X_1$ is cartesian.  \hfill{$\square$}

\end{proof}

We conclude this section with some facts all due to Grothendieck \cite{GrothFond}.

A cleavage for a fibration $P:\CX\lra \CC$ allows the definition of a pseudofunctor 
$$\CX^-:\CC^{\mathrm{op}} \lra \mathrm{Cat} \ \ .$$
The value at $U \in \CX$ is the fibre $\CX^U$ of $P$ over $U$: it is the subcategory of $\CX$ consisting of the objects $X \in \CX$ with $PU=U$ and morphisms $x:Y\lra X$ with $Px=1_U$. (This agrees with the previous notation $\CX^C$ when $C=U$ is discrete.) The value of the pseudofunctor at $h:V\lra U$ in $\CC$ is the functor 
$$h^{\star} : \CX^U \lra \CX^V$$
taking $X$ to $h^{\star}X$ and using the universal property of cartesian morphisms to define $h^{\star}$ on morphisms. If $P:\CX\lra \CC$ is a split fibration, $\CX^-:\CC^{\mathrm{op}} \lra \mathrm{Cat}$ is a functor. If $\mathrm{Spl}_{/\CC}$ is the 2-category of split fibrations over $\CC$ then we obtain an equivalence of 2-categories
\begin{equation}
\mathrm{Spl}_{/\CC} \simeq [\CC^{\mathrm{op}}, \mathrm{Cat}]
\end{equation}
taking $P:\CX\lra \CC$ to $\CX^-$.


\section{Internal full subcategories}\label{IFSC}

Suppose $\CC$ is a category admitting pullbacks. Let $\CC^{\mathbf{2}} =[\mathbf{2},\CC]$ be the category of morphisms in $\CC$. The functor
$$\mathrm{cod} : \CC^{\mathbf{2}} \lra \CC \ \ ,$$
taking each morphism to its codomain, is a fibration. The cartesian morphisms in $\CC^{\mathbf{2}} $ are the pullback squares in $\CC$.

For any category $C$ in $\CC$, the category $(\CC^{\mathbf{2}})^C$ has objects those functors $p:E\lra C$ in $\CC$ such that the square
\begin{equation}\label{discopfib}
\xymatrix{
E_1 \ar[rr]^-{d_0} \ar[d]_-{p_1} && E_0 \ar[d]^-{p_0} \\
C_1 \ar[rr]_-{d_0} && C_0}
\end{equation}
is a pullback. Such functors $p$ are called \textit{discrete opfibrations over} $C$. We think of them as internalized functors $C\lra \CC$. By Theorem~\ref{gYl}, each discrete opfibration $p:E\lra C$ corresponds to a cartesian-morphism-preserving functor 
\begin{equation}\label{pbptotal}
\bar{p} : \begingroup\textstyle\int\endgroup C \lra \CC^{\mathbf{2}}
\end{equation}
taking $(U,u)$ to $p_u:E_u\lra U$ where the square 
\begin{equation}\label{fibreoveru}
\xymatrix{
E_u \ar[rr]^-{i_u} \ar[d]_-{p_u} && E \ar[d]^-{p} \\
U \ar[rr]_-{u} && C}
\end{equation}
is a pullback. For a morphism $(h,\theta):(U,u)\lra (V,v)$, we obtain the morphism $\bar{p}(h,\theta)$:
\begin{equation}
\xymatrix{
E_u \ar[rr]^-{\bar{\theta}} \ar[d]_-{p_u} && E_v \ar[d]^-{p_v} \\
U \ar[rr]_-{h} && V}
\end{equation}
in $\CC^{\mathbf{2}}$ using the pullback \eqref{fibreoveru} for $v$ and the diagram
$$
\xymatrix{
E_u \ar[rr]^-{d_1(\theta p_u,p_0i_u)} \ar[d]_-{p_u} &  & E \ar[d]^-{p} \\
U \ar[r]_-{h} & V \ar[r]_-{v} & C }
$$
for which the pullback \eqref{discopfib} is used to define $(\theta p_u,p_0i_u):E_u\lra E_1$ satisfying $$d_0(\theta p_u,p_0i_u) =i_u \ \ \mathrm{and} \ \ p_1(\theta p_u,p_0i_u) = \theta p_u \ \ .$$ 

The split fibration $P:\int{C} \lra \CC$ corresponds to the functor 
$$\CC(-,C):\CC^{\mathrm{op}}\lra \mathrm{Cat} \ \ .$$
A choice of pullbacks in $\CC$ means that the fibration $\mathrm{cod} : \CC^{\mathbf{2}} \lra \CC$ leads to a pseudofunctor $$\CC/-: \CC^{\mathrm{op}}\lra \mathrm{Cat}$$ taking each object $U$ to the slice category $\CC/U$ and defined on morphisms by pullback along them. It follows that $\bar{p}$ induces a pseudonatural transformation with component at $U\in \CC$ denoted by 
\begin{equation}\label{pbp}
\bar{p}^U:\CC(U,C)\lra \CC/U \ \ .
\end{equation}
\begin{proposition}\label{ffonfib} A functor 
\begin{equation}\label{funover}
\xymatrix{
\CX \ar[rd]_{P}\ar[rr]^{F}   && \CY \ar[ld]^{Q} \\
& \CC  &
}
\end{equation}
over $\CC$ induces a family of functors
\begin{equation} F^U:\CX^U\lra \CY^U \ \ , \ \ U \in \CC \ \ ,
\end{equation}
on the fibres. If $F$ is fully faithful then so are all the functors $F^U$, $U\in \CC$. The converse holds if $P$ is a fibration and $F$ preserves cartesian morphisms.
\end{proposition}

\begin{proof} If $y:FX\lra FX^{\prime}$ in $\CY$ is in $\CY^U$ then any $x:X\lra X^{\prime}$ in $\CX$ with $Fx=y$ is in $\CX^U$. This proves the second sentence of the Proposition. For the converse, take any $h:U\lra U^{\prime}$ in $\CC$ and $x:Z\lra X^{\prime}$ cartesian with $Px=h$. We have a diagram of four pullbacks as below.
$$
\xymatrix{
\CX^U(X,Z) \ar[rr]^-{\mathrm{incl}} \ar[d]_-{F^U} && \CX(X,Z) \ar[rr]^-{\CX(1,x)} \ar[d]^-{F} && \CX(X,X^{\prime}) \ar[d]^-{F} \\
\CY^U(FX,FZ) \ar[rr]_-{\mathrm{incl}} \ar[d]_-{!} && \CY (FX,FZ) \ar[rr]_-{\CY(1,Fx)} \ar[d]^-{Q} && \CY(FX,FX^{\prime}) \ar[d]^-{Q} \\
1 \ar[rr]_-{\ulcorner 1_U \urcorner} && \CC(U,U) \ar[rr]_-{\CC(1,h)} && \CC(U,U^{\prime})}
$$
By assumption, the top left vertical function is invertible. So $F$ is fully faithful on those morphisms over $h$; but all morphisms are over some $h$. So $F$ is fully faithful. \hfill{$\square$}

\end{proof}

An \textit{internal full subcategory} of the category $\CC$ is a discrete opfibration $p:E\lra C$ for which each of the functors $\bar{p}^U$ \eqref{pbp} is fully faithful. By Proposition~\ref{ffonfib}, this is the same as the requirement that the functor $\bar{p} : \begingroup\textstyle\int\endgroup C \lra \CC^{\mathbf{2}}$ \eqref{pbptotal} be fully faithful.

For any discrete opfibration $p:E\lra C$, we have a factorization
\begin{equation}\label{doffact}
\xymatrix{
\CC/C_0 \ar[rd]_{\mathrm{incl}}\ar[rr]^{\bar{p_0}}   && \CC^{\mathbf{2}} \\
& \begingroup\textstyle\int\endgroup C \ar[ru]_{\bar{p}}  &
}
\end{equation}
of $\bar{p_0}$ over $\CC$ in which the first functor is bijective on objects. All three functors preserve cartesian morphisms; indeed, $\CC/C_0 \lra \begingroup\textstyle\int\endgroup C$ preserves the cleavages.

For any cartesian-morphism-preserving functor $F$ \eqref{funover} over $\CC$, we can take the `full image'
\begin{equation}\label{facttri}
\xymatrix{
\CX \ar[r]^-{H} \ar[rd]_-{P} & \CZ \ar[d]^-{R} \ar[r]^-{G}  & \CY \ar[ld]^-{Q} \\
& \CC  }
\end{equation}
where $\CZ$ is the category with the same objects as $\CX$ and with homsets $\CZ(X^{\prime},X)=\CY(FX^{\prime},FX)$. The functor $H$ is the identity on objects and defined by 
$$F: \CX(X^{\prime},X)\lra \CY(FX^{\prime},FX)=\CZ(X^{\prime},X)$$
on morphisms. The functor $G$ is defined to be $F$ on objects and the identity on morphisms. The functor $R$ is unique making \eqref{facttri} commutative. The cartesian morphisms in $\CZ$ are those which are cartesian in $\CY$. If $P$ is a fibration then so is $R$; given $h:V\lra U$ in $\CC$ and $X\in \CZ$ with $RX=PX=U$, there is a cartesian $x:X^{\prime}\lra X$ over $h$ for $P$, and $Fx:X^{\prime}\lra X$ in $\CZ$ is cartesian over $h$ for $R$. Both $H$ and $G$ preserve cartesian morphisms.

Let us write $P_F : \CX[F]\lra \CC$ for the functor $R:\CZ \lra \CC$ obtained using the above full image construction.

A morphism $k:M\lra N$ in $\CC$ (can be regarded as a discrete opfibration between discrete categories in $\CC$ and so) gives rise to a cartesian-morphism-preserving functor
\begin{equation}
\xymatrix{
\CC/N \ar[rd]_{\mathrm{dom}}\ar[rr]^{\bar{k}}   && \CC^{\mathbf{2}} \ar[ld]^{\mathrm{cod}} \\
& \CC  &
}
\end{equation}
over $\CC$. We can factor $\bar{k}$ as
\begin{equation}\label{factkbar}
\xymatrix{
\CC/N \ar[r]^-{H} \ar[rd]_-{\mathrm{dom}} & (\CC/N)[\bar{k}] \ar[d]^-{P_{\bar{k}}} \ar[r]^-{ G}  & \CC^{\mathbf{2}} \ar[ld]^-{\mathrm{cod}} \\
& \CC  }
\end{equation}
where $H$ is the identity on objects and $G$ is fully faithful. The objects of $(\CC/N)[\bar{k}]$ are pairs $(U,u)$ for $u:U\lra N$ in $\CC$ and the morphisms are commutative squares
\begin{equation}
\xymatrix{
M_u \ar[rr]^-{} \ar[d]_-{k_u} && M_v \ar[d]^-{k_v} \\
U \ar[rr]_-{h} && V}
\end{equation}
where $k_u$ is the pullback of $k$ along $u$.

The following result is due to B\'enabou \cite{Ben1973Top}. It is Theorem 6.2 of \cite{CoIC} and Example 2.38 of \cite{Johnstone1977}. 
Also see \cite{Ben1975Fib}, \cite{Wraith1975} and \cite{JohnstWraith}.  

\begin{proposition}  For a morphism $k:M\lra N$ in a finitely complete category $\CC$, the following three conditions are equivalent:
\begin{itemize}
\item[(i)] there exists an internal full subcategory $p:E\lra C$ of $\CC$ with $p_0=k$;
\item[(ii)] there exists a category $C$ in $\CC$ with $C_0=N$ and an isomorphism
$${\begingroup\textstyle\int\endgroup}C\cong (\CC/N)[\bar{k}]$$ 
of categories over $\CC$;
\item[(iii)] there exists a cartesian internal hom for the two objects
$$k\times 1: M\times N \lra N\times N \ \ \mathit{and} \ \ 1\times k: N \times M \lra N\times N$$
of the category $\CC/N\times N$.
\end{itemize}
\end{proposition}

\begin{proof} (i) $\Rightarrow$ (ii) Given (i), compare the factorizations \eqref{doffact} and \eqref{factkbar}. Since bijective-on-objects and fully faithful functors form a factorization system on $\mathrm{Cat}$, the two factorizations are isomorphic, yielding (ii).

\noindent (ii) $\Rightarrow$ (iii) The product of $(u,v):U\lra N \times N$ and $k\times 1 :M\times N\lra N\times N$ in $\CC/N\times N$ is the main diagonal of the pullback
$$
\xymatrix{
M_u \ar[rr]^-{} \ar[d]_-{k_u} && M\times N \ar[d]^-{k\times 1} \\
U \ar[rr]_-{(u,v)} && N\times N}
$$
in $\CC$. So   
\begin{equation}\label{slicehom1}
(\CC/N\times N)(U\operatorname*{\times }\limits_{N\times N} (M \times N)\to N \times N,  N \times M \stackrel{1\times k}{\lra} N \times N)
\cong
(\CC/U)(M_u \stackrel{k_u}{\lra} U,M_v \stackrel{k_v}{\lra} U) \ .
\end{equation}
On the other hand, for any category $C$ with $C_0=N$, we have
\begin{equation}\label{slicehom2}
(\CC/N\times N)(U\stackrel{(u,v)}{\lra} N\times N, C_1 \stackrel{(d_0,d_1)}{\lra} N\times N)\cong \CC(U,C)(u,v) \ .
\end{equation}
The isomorphism of (ii) yields an isomorphism between the right hand side of \eqref{slicehom1} and the right hand side of \eqref{slicehom2}. This shows that 
$$(d_0,d_1):C_1\lra N\times N$$
is an internal hom as required for (iii).

\noindent (iii) $\Rightarrow$ (i) Let $(d_0,d_1):C_1\lra N\times N$ be an internal hom as in (iii) and let $C_0=N$. Despite $C$ not yet being a category, only a graph, the left hand side of \eqref{slicehom1} and that of \eqref{slicehom2} are isomorphic. Therefore, looking at the right hand sides, we have a fully faithful graph morphism
$$\CC(U,C) \lra \CC/U \ , \ \ u \mapsto k_u \ .$$ 
Since the codomain is a category, a category structure is induced on $\CC(U,C)$. The family of functors is pseudonatural in $U$ and so yields a fully faithful cartesian-morphism-preserving functor $\int C \lra \CC^{\mathbf{2}}$ over $\CC$. By Theorem~\ref{gYl}, (i) follows. \hfill{$\square$}
\end{proof}

\section{Slices and left extensions}\label{Sle} 

For functors $F:\CA \lra \CC$ and $G:\CB \lra \CC$, the so-called \textit{comma category} $F/G$ (for example, see \cite{CWM}) has objects $(A,h:FA\lra GB,B)$ where $A$ is an object of $\CA$, $B$ is an object of $\CB$, and $h:FA\lra GB$ is a morphism of $\CC$. A morphism $(f,g): (A,h,B)\lra (A^{\prime},h^{\prime},B^{\prime})$ in $F/G$ consists of a morphism $f:A\lra A^{\prime}$ in $\CA$ and a morphism $f:B\lra B^{\prime}$ in $\CB$ such that $h^{\prime}\circ Ff = Gg \circ h$. We write $P:F/G \lra \CA$ and $Q:F/G\lra \CB$ for the projection functors. The special case where $F$ is the identity functor of $\CC$ and $\CB$ is the terminal category, so that $G$ can be identified with an object $U$ of $\CC$, is called traditionally called the \textit{slice} of $\CC$ over $U$.  

We propose to use the term \textit{slice of} $F$ \textit{over} $G$ for the general comma category $F/G$. This terminology seems to work better than my traditional term `comma object' \cite{FYl} for the construction internalized to a pair of morphisms in a 2-category (recalled below). Weber \cite{Weber2007} uses the term `lax pullback'.

Before internalizing definitions to a 2-category, we remind the reader (see \cite{CWM}) 
that the left Kan extension $K:\CB \lra \CC$ of a functor $F:\CA \lra \CC$ along a functor $H:\CA \lra \CB$ can be calculated by Lawvere's pointwise formula
$$KB = \mathrm{colim} (H/B \stackrel{P}\lra \CA \stackrel{F}\lra \CC)$$
when the colimits exist. Furthermore, the colimit of a functor is none other than a left Kan extension along a functor into the terminal category $\mathbf{1}$. 

Let $\CK$ be a 2-category. In $\CK$, a \textit{slice} of a morphism $f:A\lra C$ over a morphism $g:B\lra C$ is a square
\begin{equation}\label{slicesq}
\xymatrix{
f/g \ar[d]_{p}^(0.5){\phantom{aaa}}="1" \ar[rr]^{q}  && B \ar[d]^{g}_(0.5){\phantom{aaa}}="2" \ar@{=>}"1";"2"^-{\lambda}
\\
A \ar[rr]_-{f} && C 
}
\end{equation} 
such that, for all objects $X$, the functor 
\begin{equation}\label{slicereptble} 
\CK(X,f/g) \lra \CK(X,f)/\CK(X,g) \ ,
\end{equation} 
taking $r:X\lra f/g$ to $(pr,\lambda r, qr)$, is an isomorphism. A slice of the identity morphism $1_C:C\lra C$ over itself is the \textit{cotensor} $C^{\mathbf{2}}$ \textit{of} $C$ \textit{with the arrow category} $\mathbf{2}$.   

A 2-category $\CK$ is \textit{(finitely) complete} \cite{Licv2f} when it has (finite) products and equalizers (in the $\mathrm{Cat}$-enriched sense), and has cotensors with $\mathbf{2}$. It follows that all slices exist. Also all cotensors with (finite) categories exist. 

If $\CC$ is a (finitely) complete category then $\CK=\mathrm{Cat}\CC$ is a (finitely) complete 2-category; see \eqref{arrobj} for the construction of cotensors with $\mathbf{2}$. By Corollary~\ref{discretedetectlim}, the square \eqref{slicesq} is a slice if and only if the functor \eqref{slicereptble} is invertible for all discrete $X$.    

The 2-category $\mathrm{Cat}/\CC$ is complete. In particular, we are interested in the slice $F/_{\CC}F^{\prime}$ in $\mathrm{Cat}/\CC$ of two functors over $\CC$ as in the following diagram.
\begin{equation}
\xymatrix{
\CX \ar[r]^-{F} \ar[rd]_-{P} & \CY \ar[d]^-{Q}  & \CX^{\prime} \ar[l]_-{F^{\prime}} \ar[ld]^-{P^{\prime}} \\
& \CC  }
\end{equation}
Indeed, $F/_{\CC}F^{\prime}$ is the full subcategory of the slice category $F/F^{\prime}$ consisting of those objects $(X,y,X^{\prime})$, where $X\in \CX$, $X^{\prime}\in \CX^{\prime}$ and $y\in \CY(FX,F^{\prime}X^{\prime})$, for which $Qy$ is an identity morphism in $\CC$. We have a universal square containing a natural transformation, as in the following square, taken to an identity by $Q:\CY\lra \CC$.
\begin{equation}
\xymatrix{
F/_{\CC}F^{\prime} \ar[d]_{\mathrm{dom}}^(0.5){\phantom{aaaa}}="1" \ar[rr]^{\mathrm{cod}}  && \CX^{\prime} \ar[d]^{F^{\prime}}_(0.5){\phantom{aaaa}}="2" \ar@{=>}"1";"2"^-{\lambda}
\\
\CX \ar[rr]_-{F} && \CY 
}
\end{equation}
If $F$ and $F^{\prime}$ preserve cartesian morphisms then so do $\mathrm{dom}$ and $\mathrm{cod}$, and the cartesian morphisms of $F/_{\CC}F^{\prime}$ are precisely those taken to cartesian morphisms by $\mathrm{dom}$ and $\mathrm{cod}$. 

\begin{proposition} The 2-functor 
$$\begingroup\textstyle\int\endgroup : \mathrm{Cat}\CC\lra \mathrm{Cat}/\CC$$
preserves slices.
\end{proposition}

\begin{proof} The universal property of a slice
\begin{equation}
\xymatrix{
f/f^{\prime} \ar[d]_{d}^(0.5){\phantom{aaaa}}="1" \ar[rr]^{d^{\prime}}  && C^{\prime} \ar[d]^{f^{\prime}}_(0.5){\phantom{aaaa}}="2" \ar@{=>}"1";"2"^-{\lambda}
\\
C \ar[rr]_-{f} && D 
}
\end{equation}
in $\mathrm{Cat}\CC$ immediately yields that the functor 
$$\begingroup\textstyle\int\endgroup (f/f^{\prime}) \lra {\begingroup\textstyle\int\endgroup} f /_{\CC} {\begingroup\textstyle\int\endgroup} f^{\prime} \ \  , $$
taking the morphism
$$
\xymatrix{
U \ar[rd]_{u}^(0.5){\phantom{A}}="1" \ar[rr]^{h}  && V \ar[ld]^{v}_(0.5){\phantom{A}}="2" \ar@{=>}"1";"2"^-{\theta}
\\
& f/f^{\prime} 
}
$$
in ${\begingroup\textstyle\int\endgroup}f/_{\CC}{\begingroup\textstyle\int\endgroup}f^{\prime}$ to the morphism
$$((h,d\theta),(h,d^{\prime}\theta)):((U,du),\lambda u,(U,d^{\prime} u))\lra((V,dv),\lambda v,(V,d^{\prime} v))$$
in ${\begingroup\textstyle\int\endgroup}f/_{\CC}{\begingroup\textstyle\int\endgroup}f^{\prime}$, is an isomorphism.    \hfill{$\square$}
\end{proof}

In a 2-category $\CK$, a diagram
\begin{equation}\label{lextn}
\xymatrix{
A \ar[rd]_{f}^(0.5){\phantom{a}}="1" \ar[rr]^{h}  && B \ar[ld]^{k}_(0.5){\phantom{a}}="2" \ar@{=>}"1";"2"^-{\kappa}
\\
& C 
}
\end{equation}
is said \cite{FTM} to exhibit $k$ as a \textit{left extension} of 
$f$ along $h$ when, for all $g:B\lra C$, pasting with the triangle 
yields a bijection between 2-cells $k\Lra g$ 
and 2-cells $f\Lra gh$. To say such a left extension exists for 
each $f$, is to say the functor 
$$\CK(h,C) : \CK(B,C) \lra \CK(A,C)$$
has a left adjoint.

The diagram \eqref{lextn} is said \cite{FYl} to exhibit $k$ as 
a \textit{pointwise left extension} of $f$ along $h$ when, for all
morphisms $b:X\lra B$ for which the slice of $h$ over $b$ exists,
the diagram
\begin{equation}\label{pwlextn}
\xymatrix{h/b \ar[d]_{p}^(0.5){\phantom{aaaaaa}}="1" \ar[rr]^{q}  && X \ar[d]^{b}_(0.5){\phantom{aaaaaa}}="2" \ar@{=>}"1";"2"^-{\lambda}
\\
A \ar[rd]_{f}^(0.5){\phantom{aa}}="1" \ar[rr]^{h}  && B \ar[ld]^{k}_(0.5){\phantom{aa}}="2" \ar@{=>}"1";"2"^-{\kappa}
\\
& C 
}
\end{equation}
exhibits $kb$ as a left extension of $fp$ along $q$. It is shown in \cite{FYl} that every pointwise left extension is a left extension, provided the slice $h/1_B$ exists, and that
\eqref{pwlextn} furthermore exhibits $kb$ as a pointwise left extension of $fp$ along $q$. 
  
\begin{definition}\label{ff} \em{\cite{FYl} A morphism $f:A\lra B$ in $\CK$ is \textit{(representably) fully faithful} when the functor $\CK (X,f):\CK (X,A) \lra \CK (X,B)$ is fully faithful for all objects $X$ of $\CK$. This is equivalent, when the cotensor with $\mathbf{2}$ exists, to saying that the square
$$
\xymatrix{
A^{\mathbf{2}} \ar[d]_{d_0}^(0.5){\phantom{aaaa}}="1" \ar[rr]^{d_1}  && A \ar[d]^{f}_(0.5){\phantom{aaaa}}="2" \ar@{=>}"1";"2"^-{f\lambda}
\\
A \ar[rr]_-{f} && B 
}
$$
exhibits $A^{\mathbf{2}}$ as the slice $f/f$.}
\end{definition}

\begin{definition}\label{dense} \em{\cite{EC} A morphism $f:A\lra B$ in $\CK$ is {\em{dense}}  when the triangle
$$
\xymatrix{
A \ar[rd]_{f}^(0.5){\phantom{a}}="1" \ar[rr]^{f}  && B \ar[ld]^{1_B}_(0.5){\phantom{a}}="2" \ar@{=>}"1";"2"^-{1_f}
\\
& B 
}
$$
exhibits $1_B$ as a pointwise left extension of $f$ along $f$.}
\end{definition}

\section{Descent categories}\label{Dc}

As usual we write $\mathbf{\Delta}$ for the (topologists') simplicial category. The objects are the non-empty linearly ordered sets 
$$[n]=\{0,1,\dots,n\}$$ 
for $n\geq 0$, and the morphisms are order-preserving functions. We write
$$\partial_p:[n-1]\lra [n] \ \  \text{ and } \ \  \sigma_q:[n+1]\lra [n]$$
for the injective order-preserving function which omits the element $p$ in its image and the surjective order-preserving function which identifies only the elements $q$ and $q+1$. The squares
$$
\xymatrix{
[n-1] \ar[rr]^-{\partial_0} \ar[d]_-{\partial_n} && [n] \ar[d]^-{\partial_{n+1}} \\
[n] \ar[rr]_-{\partial_0} && [n+1]}
$$
are pushouts in $\mathbf{\Delta}$.
 
\noindent {\bf Convention} When thinking of $\partial_p$ and $\sigma_q$ as morphisms of the category $\mathbf{\Delta}^{\mathrm{op}}$, we write them respectively as
$$d_p:[n]\lra [n-1]  \ \  \text{ and } \ \  i_q:[n] \lra [n+1] \ .$$

Let $\mathbf{D}$ denote the category
\begin{equation}\label{D}
\xymatrix{
[3] \ar@<+3ex>[rr]|-{d_{0}} \ar@<+1ex>[rr]|-{d_{1}} \ar@<-1ex>[rr]|-{d_{2}} \ar@<-3ex>[rr]|-{d_{3}} && \ar@<+4ex>[rr]|-{d_{0}} \ar@{<-}@<+2ex>[rr]|-{i_{0}} \ar[rr]|-{d_{1}} \ar@{<-}@<-2ex>[rr]|-{i_{1}} \ar@<-4ex>[rr]|-{d_{2}} [2] && [1] \ar@<+2ex>[rr]|-{d_{0}} \ar@{<-}[rr]|-{i} \ar@<-2ex>[rr]|-{d_{1}} && [0]
}
\end{equation}
in $\mathbf{\Delta}^{\mathrm{op}}$.

The origins of the following construction go back to Grothendieck \cite{GrothDesc}.

\begin{definition}\label{Desc} \em{The \textit{descent category} for a functor $P:\CE \lra \mathbf{\Delta}^{\mathrm{op}}$ is 
$$\mathrm{Desc}\CE = \CE^{\mathbf{D}} \ .$$
 Objects $E$ of $\mathrm{Desc}\CE$ are called \textit{descent data} for $P$}.
\end{definition} 

If, in Definition~\ref{Desc}, the functor $P$ is a fibration, a cleavage determines a pseudofunctor
$$\CE^{-} : \mathbf{\Delta} \lra \mathrm{Cat} \  ;$$
that is, a pseudocosimplicial category. We write $\CE_n$ for the fibre $\CE^{[n]}$ of $P$ over $[n]$, yielding a diagram
\begin{equation}
\xymatrix{
\CE_{3} \ar@{<-}@<+3ex>[rr]|-{\partial_{0}} \ar@{<-}@<+1ex>[rr]|-{\partial_{1}} \ar@{<-}@<-1ex>[rr]|-{\partial_{2}} \ar@{<-}@<-3ex>[rr]|-{\partial_{3}} && \ar@{<-}@<+4ex>[rr]|-{\partial_{0}} \ar@<+2ex>[rr]|-{\sigma_{0}} \ar@{<-}[rr]|-{\partial_{1}} \ar@<-2ex>[rr]|-{\sigma_{1}} \ar@{<-}@<-4ex>[rr]|-{\partial_{2}} \CE_{2} && \CE_{1} \ar@{<-}@<+2ex>[rr]|-{\partial_{0}} \ar [rr]|-{\sigma} \ar@{<-}@<-2ex>[rr]|-{\partial_{1}} && \CE_{0}
}
\end{equation}
in $\mathrm{Cat}$ in which the cosimplicial identities hold up to coherent isomorphism: for example,
$$\partial_q \partial_p \cong \partial_p \partial_{q-1} \text{ \  for \ } p<q \ .$$

For descent data $E$, define
\begin{equation}\label{eDesc}
e:\partial_0 E_0 \lra \partial_1 E_0
\end{equation}
to be the unique morphism in $\CE_1$ such that the triangle
$$\xymatrix{
E_1 = \partial_0 E_0 \ar[rd]_{d_1}\ar[rr]^{e}   && \partial_1 E_0 \ar[ld]^{{d_1}^{E_0}} \\
&  E_0  &
}
$$
commutes. We readily see that the hexagon
\begin{equation}\label{ehexDesc}
\xymatrix{
& \partial_1\partial_0E_0 \ar[r]^-{\partial_1e} & \partial_1\partial_1E_0 \ar[rd]^-{\cong} & \\
\partial_0\partial_0E_0 \ar[ru]^-{\cong} \ar[rd]_-{\partial_0e} & & & \partial_2\partial_1E_0 \\
& \partial_0\partial_1E_0 \ar[r]_-{\cong} & \partial_2\partial_0E_0 \ar[ru]_-{\partial_2e} &}
\end{equation}
commutes in $\CE_2$ and
\begin{equation}\label{etriDesc}
\xymatrix{
\sigma_0\partial_0E_0 \ar[rd]_{\cong}\ar[rr]^{\sigma_0e}   && \sigma_0\partial_1E_0 \ar[ld]^{\cong} \\
& E_0  &
}
\end{equation}
commutes in $\CE_0$. In fact, we can reconstruct $E$ from $E_0$ and $e$ \eqref{eDesc} satisfying \eqref{ehexDesc} and \eqref{etriDesc}.

While we have defined the descent category, for a category `parametrized' by $\mathbf{\Delta}^{\mathrm{op}}$, in terms of the $\CX^C$ construction, we can go the other way too. Suppose $P:\CX \lra \CC$ is a functor and $C$ is a category in $\CC$. For simplicity we suppose $\CC$ has pullbacks chosen so that we can complete the diagram for $C$ to a simplicial object (functor)
$$C:\mathbf{\Delta}^{\mathrm{op}} \lra \CC \ .$$
Form the pullback
\begin{equation}\label{PCpb}
\xymatrix{
C^{\star}\CX\ar[rr]^-{ } \ar[d]_-{P_C} && \CX \ar[d]^-{P} \\
\mathbf{\Delta}^{\mathrm{op}} \ar[rr]_-{C} && \CC}
\end{equation}
in $\mathrm{Cat}$.
\begin{proposition} There is a canonical equivalence of categories 
$$\CX^C \simeq \mathrm{Desc}C^{\star}\CX$$
\[
\xymatrix @C+5mm {
E \ar @{<->} [r] |-*@{|} & (\mathbf{D},E) }.
\]
For cloven $P$, this is an isomorphism.
\end{proposition}

For any functor $P:\CX\lra \CC$ and category $\CH$, the cotensor of $\CH$ with $P$ in the 2-category $\mathrm{Cat}/\CC$ is defined by the pullback
$$
\xymatrix{
[\CH,\CX]_{\CC} \ar[rr]^-{ } \ar[d]_-{P^{\CH}} && [\CH,\CX] \ar[d]^-{[1,P]} \\
\CC \ar[rr]_-{\mathrm{diag}} && [\CH,\CC]}
$$
in $\mathrm{Cat}$. For a cloven fibration $P$, the pseudofunctor corresponding to the cloven fibration $P^{\CH}$ is the composite
$$\CC^{\mathrm{op}} \stackrel{\CX^{-}} {\lra} \mathrm{Cat} \stackrel{[\CH,-]} {\lra} \mathrm{Cat} \ .$$ 
\begin{proposition}\label{DescRep} For any functor $P:\CE \lra \mathbf{\Delta}^{\mathbf{op}}$ and category $\CH$, there is a canonical isomorphism of categories
$$[\CH,\mathrm{Desc}\CE] \cong \mathrm{Desc}[\CH,\CE]_{\mathbf{\Delta}^{\mathbf{op}}} \ . 
$$
\end{proposition}

Loosely speaking, this says that the descent construction is preserved by representable 2-functors out of $\mathrm{Cat}$. It therefore makes sense to define the construction in a 2-category $\CK$.

\section{Descent and codescent objects} 

For a pseudofunctor $T:\mathbf{\Delta} \lra \mathrm{Cat}$, we write $\mathrm{Desc}T$ for the category $\mathrm{Desc}\CE$ where $P:\CE \lra \mathbf{\Delta}^{\mathbf{op}}$ is the cloven fibration obtained from $T$ by the Grothendieck construction described at the end of Section~\ref{FgYL}.

\begin{definition}\label{descobj} For a pseudofunctor 
$$T:\mathbf{\Delta}\lra \CK \ ,$$
a \textit{descent object} is an object $\mathrm{Desc}T$ of $\CK$ equipped with a 2-natural isomorphism
$$\CK(K, \mathrm{Desc}T) \cong \mathrm{Desc}\CK(K,T) \ .$$
A \textit{codescent object} in $\CK$ for a pseudofunctor $S: \mathbf{\Delta}^{\mathbf{op}} \lra \CK$ is a descent object for $S^{\mathrm{op}} : \mathbf{\Delta} \lra \CK^{\mathrm{op}}$. 
\end{definition}

\begin{proposition} A descent object for a functor $T: \mathbf{\Delta} \lra \CK$ is a limit for $T$ weighted by the inclusion $\mathbf{\Delta} \lra \mathrm{Cat}$.
\end{proposition}

B\'enabou-Roubaud \cite{BenRou1970} and Jon Beck (unpublished) expressed descent data in terms of Eilenberg-Moore algebras for monads. We shall describe a version of this. 

A functor $P:\CX \lra \CC$ is called an \textit{opfibration} (originally Grothendieck used the term `cofibration' and some authors persist with this) when $P^{\mathbf{op}}:\CX^{\mathbf{op}} \lra \CC^{\mathbf{op}}$ is a fibration. Cartesian morphisms and inverse images for $P^{\mathbf{op}}$ are called \textit{opcartesian morphisms} and \textit{direct images} for $P$.

\begin{proposition} If a cloven fibration $P:\CX \lra \CC$ is also an opfibration then the inverse image functor
$$h^{\star}:\CX^U \lra \CX^V \ ,$$
for each $h:V \lra U$ in $\CC$, has a left adjoint defined by direct image $h_{\star}$ along $h$.
\end{proposition}

\begin{proof} $\CX^V(Y,h^{\star}X) \cong \{x:Y \lra X | Px=h\} \cong \CX^U(h_{\star}Y,X) \ .$ \hfill{$\square$}
\end{proof}

For any functor $P:\CX \lra \CC$, consider a commutative square
\begin{equation}\label{commsq}
\xymatrix{
R \ar[rr]^-{q} \ar[d]_-{p} && Y \ar[d]^-{y} \\
X \ar[rr]_-{x} && Z}
\end{equation}
in $\CX$. The following condition is attributed to Chevalley in \cite{BenRou1970} and to Jon Beck in other places.

\noindent \textbf{Chevalley-Beck condition} \textit{If \eqref{commsq} is a pullback in} $\CX$ \textit{preserved by} $P$ \textit{with} $x$ \textit{cartesian and} $y$ \textit{opcartesian then} $p$ \textit{is opcartesian}.  

By Lemma~\ref{cart&pb}, the hypothesis here that \eqref{commsq} be a pullback in $\CX$ is equivalent to the assumption that $q$ be cartesian. 

\begin{example}\label{CBforcod} 
\em{The Chevalley-Beck condition holds for $\mathrm{cod}:\CC^{\mathbf{2}} \lra \CC$ for any category $\CC$. To see this, notice that a square \eqref{commsq} in $\CC^{\mathbf{2}}$ is a cube \eqref{cube}.
\begin{equation}\label{cube}
\xymatrix{
R \ar[rrr]^-{q_0} \ar[rd] \ar[ddd]_-{p_0} & & & Y \ar[ld] \ar[ddd]^-{y_0} \\
& M \ar[d]_-{p_1} \ar[r]^-{q_1} & V \ar[d]^-{y_1} \\
& U \ar[r]_-{x_1} & W \\
X \ar[ru] \ar[rrr]_-{x_0} & & & Z \ar[lu] } 
\end{equation}
The assumptions of Chevalley-Beck are that the front, back, top and bottom faces of \eqref{cube} are pullbacks, and that $y_0$ is invertible. Since the top face is a pullback, $p_0$ is also invertible. So $p$ is opcartesian.}
\end{example}

If $P$ is a cloven fibration and opfibration, the Chevalley-Beck condition is the requirement that, for all pullback squares
\begin{equation}\label{CBgamma}
\xymatrix{
M \ar[rr]^-{q} \ar[d]_-{p} && V \ar[d]^-{k} \\
U \ar[rr]_-{h} && W}
\end{equation}
in $\CC$, the canonical natural transformation
\begin{equation}
\xymatrix{
\CX^U \ar[d]_{p^{\star}}^(0.5){\phantom{aaa}}="1" \ar[rr]^{h_{\star}}  && \CX^W \ar[d]^{k^{\star}}_(0.5){\phantom{aaa}}="2" \ar@{=>}"1";"2"^-{\gamma}
\\
\CX^M \ar[rr]_-{q_{\star}} && \CX^V 
}
\end{equation}
is invertible. Under these circumstances, each category $C$ in $\CC$ determines a monad $T_C$ on the category $\CX^{C_0}$ as follows. The endofunctor for the monad is the composite
$$ \CX^{C_0} \stackrel{d_0^{\star}} {\lra}\CX^{C_1} \stackrel{d_{1\star}} {\lra}\CX^{C_0} \ . $$
The unit is the composite natural transformation
$$ 1 \cong (d_1 i)_{\star} (d_0 i)^{\star} \cong d_{1\star}i_{\star}i^{\star}d_0^{\star} 
\stackrel{d_{1\star}\varepsilon d_0^{\star}} {\lra}d_{1\star}d_0^{\star}$$
where $\varepsilon : i_{\star}i^{\star} \lra 1$ is the counit for $i_{\star} \dashv i^{\star}$. The multiplication is the composite natural transformation
$$d_{1\star}d_0^{\star}d_{1\star}d_0^{\star}\stackrel{d_{1\star}\gamma^{-1}d_0^{\star}}\lra d_{1\star}d_{2\star}d_0^{\star}d_0^{\star} \cong d_{1\star}d_{1\star}d_1^{\star}d_0^{\star}\stackrel{d_{1\star}\varepsilon d_0^{\star}} {\lra}d_{1\star}d_0^{\star}$$
using the $\gamma^{-1}$ of \ref{CBgamma} for a pullback in condition (C4) of Section \ref{Cic} and the counit $\varepsilon$ for $d_{1\star} \dashv d_0^{\star}$.

The following result appears as Proposition (9.10) of \cite{CoIC}.

\begin{proposition}\label{FibasEM} 
If $P:\CX \lra \CC$ is a cloven fibration and opfibration satisfying 
the Chevalley-Beck condition then, for each category $C$ in $\CC$ 
there is an equivalence of categories 
$$\CX^C \simeq (\CX^{C_0})^{T_C}$$ 
where the right-hand side is the category of Eilenberg-Moore algebras for the monad $T_C$.
\end{proposition}

\begin{proof} Just as for descent data with \eqref{eDesc}, for each $X\in\CX^C$, we obtain a unique morphism
\begin{equation}
e:d_0^{\star} X_0 \lra d_1^{\star} X_0
\end{equation}  
in $\CX^C_1$ defined by the equation 
$$(d_0^{\star} X_0  \stackrel{e} {\lra} d_1^{\star} X_0 \stackrel{d_1^{X_0}}{\lra}X_0) = (d_0^{\star} X_0 = X_1 \stackrel{d_1}\lra X_0) \ .$$
Diagrams as in \eqref{ehexDesc} and \eqref{etriDesc} hold, yielding that the mate 
$$\hat{e}:T_C X_0=d_{1\star}d_0^{\star} X_0 \lra X_0 \ ,$$
of $e$ under the adjunction $d_{1\star}\dashv d_1^{\star}$, is a $T_C$-algebra structure on $X_0$. The remainder is left to the reader. \hfill{$\square$}
\end{proof}

\begin{example} \em{For the moment, consider $\mathrm{Cat}$ as a category and let us look at the fibration 
$$\mathrm{cod}: \mathrm{Cat}^{\mathbf{2}} \lra \mathrm{Cat} \ .$$
The cartesian morphisms for a cleavage are the chosen pullbacks in $\mathrm{Cat}$. The chosen opcartesian morphisms over $H:\CC \lra \CD$ are commutative squares of the following form.
\[
\xymatrix{
\CX \ar[rr]^-{1_{\CX}} \ar[d]_-{} && \CX \ar[d]^-{} \\
\CC \ar[rr]_-{H} && \CD}
\]
As we saw using \eqref{cube}, the Chevalley-Beck condition holds. For any category $\CC$, we have a category

\begin{equation}\label{sqdoubcat}
\xymatrix{
& \mathrm{sq}\CC :  \CC^{\mathbf{3}} \ar@<+2ex>[rr]|-{ } \ar@<+0ex>[rr]|-{\mathrm{comp}} \ar@<-2 ex>[rr]|-{ } &&   \CC^{\mathbf{2}} \ar@<+2ex>[rr]|-{\mathrm{dom}} \ar@{<-}[rr]|-{\mathrm{id}} \ar@<-2ex>[rr]|-{\mathrm{cod}} && \CC 
}
\end{equation}
in $\mathrm{Cat}$ obtained as
$$\mathbf{\Delta}^{\mathrm{op}} \stackrel{\mathrm{incl}}\lra \mathrm{Cat}^{\mathrm{op}} \stackrel{[-,\CC]} \lra \mathrm{Cat} \ .$$
The reason for the name $\mathrm{sq}\CC$ is because it is the `double category of squares' in the category $\CC$. Proposition \ref{FibasEM} yields an equivalence of categories
\begin{equation}\label{sqfun}
(\mathrm{Cat})^{\mathrm{sq}\CC} \simeq (\mathrm{Cat}/\CC)^{T_{\mathrm{sq}\CC}} \ .
\end{equation}
The monad $T_{\mathrm{sq}\CC}$ is easily described. 
The endofunctor of $\mathrm{Cat}/\CC$ takes a functor $P:\CX\lra \CC$ to the split opfibration 
$$\mathrm{cod} : P/\CC \lra \CC$$
with domain the slice $P/\CC$ of the functor $P:\CX\lra \CC$ over $1_{\CC} : \CC \lra \CC$. 
The objects of $P/\CC$ are $(X,h,U)$ where $X\in \CX$ and $h:PX\lra U$ in $\CC$. The multiplication of the monad at $P$ is the functor
\[
\xymatrix{
\mathrm{cod}/\CC \ar[rd]_{\mathrm{cod}}\ar[rr]^{\mu}   && P/\CC \ar[ld]^{\mathrm{cod}} \\
& \CC  &}
\]
taking $(X,h,U,r,V)$ to $(X,rh,V)$ and the unit of the monad at $P$ is the functor
\[
\xymatrix{
\CX \ar[rd]_{P}\ar[rr]^{\eta}   && P/\CC \ar[ld]^{\mathrm{cod}} \\
& \CC  &}
\]
taking $X$ to $(X,1_{PX},PX)$. In fact, the categories in \eqref{sqfun} 
are 2-categories in an obvious way ($T_{\mathrm{sq}\CC}$ 
is a 2-monad on $\mathrm{Cat}/\CC$) and the equivalence is an equivalence of 2-categories. 
The $T_{\mathrm{sq}\CC}$-algebras are easily identified as split opfibrations $P:\CX \lra \CC$; 
the $T_{\mathrm{sq}\CC}$-action is defined by the functor 
$$\alpha : P/\CC \lra \CX$$
taking $(X,h,U)$ to $h_{\star}X \in \CX^U$. If we write $\mathrm{Spl}_{\mathrm{op}\CC}$ for
the 2-category of split opfibrations over $\CC$, cleavage preserving functors over $\CC$, 
and natural transformations over $\CC$, we see that \eqref{sqfun} can be prolonged to equivalences of 2-categories
\begin{equation}
(\mathrm{Cat})^{\mathrm{sq}\CC} \simeq (\mathrm{Cat}/\CC)^{T_{\mathrm{sq}\CC}} \cong \mathrm{Spl}_{\mathrm{op}\CC} \simeq [\CC,\mathrm{Cat}] \ .
\end{equation}}
\end{example}

\begin{example}\label{codescofdiscrete}
\em{For any category $B$ in $\CC$, the codescent object of the simplicial object
\begin{equation}\label{nerveB}
\xymatrix{
& B_2 \ar@<+2ex>[rr]|-{d_0} \ar@<+0ex>[rr]|-{d_1} \ar@<-2 ex>[rr]|-{d_2} &&   B_1 \ar@<+2ex>[rr]|-{d_0} \ar@{<-}[rr]|-{i} \ar@<-2ex>[rr]|-{d_1} && B_0 }
\end{equation}
in $\mathrm{Cat}\CC$ is easily seen to be $B$ itself. For, the definitions at the beginning of
Section \ref{Cic} make it clear that 
$$(\mathrm{Cat}\CC)(B,C)$$
is the descent category for the simplicial category
\begin{equation}
\xymatrix
{\dots   \ar@{<-}@<+4ex>[rr]|-{\partial_{0}} \ar@<+2ex>[rr]|-{\sigma_{0}} \ar@{<-}[rr]|-{\partial_{1}} \ar@<-2ex>[rr]|-{\sigma_{1}} \ar@{<-}@<-4ex>[rr]|-{\partial_{2}} (\mathrm{Cat}\CC)(B_2,C) && (\mathrm{Cat}\CC)(B_1,C) \ar@{<-}@<+2ex>[rr]|-{\partial_{0}} \ar [rr]|-{\sigma} \ar@{<-}@<-2ex>[rr]|-{\partial_{1}} && (\mathrm{Cat}\CC)(B_0,C)
}
\end{equation}
where $\partial_p = (\mathrm{Cat}\CC) (d_p,1_C)$. Observe too that, for each $U$ in $\CC$, the nerve of the category $\CC(U,B)$ is
\[
\xymatrix{
\dots \CC(U,B_2) \ar@<+2ex>[rr]|-{\CC(1,d_0)} \ar@<+0ex>[rr]|-{\CC(1,d_1)} \ar@<-2 ex>[rr]|-{\CC(1,d_2)} &&    \CC(U,B_1) \ar@<+2ex>[rr]|-{ \CC(1,d_0)} \ar@{<-}[rr]|-{\CC(1,i)} \ar@<-2ex>[rr]|-{ \CC(1,d_1)} &&  \CC(U,B_0) \ ; }
\]
so the codescent object of \eqref{nerveB} is preserved by the 2-functors 
$\CC(U,-) : \mathrm{Cat}\CC \lra \mathrm{Cat}$. 
This provides another proof of Corollary \ref{discretedensity} that $\CC \lra \mathrm{Cat}\CC$ is dense, by showing that there is a density presentation via a codescent construction. }
\end{example}

\section{Split fibrations in a category}\label{SFiaC}

Let $\CC$ be a category with pullbacks. For each category $B$ in $\CC$, we can construct a category $B^{\mathbf{2}}$ in $\CC$ as follows. Form the diagram 
$$
\xymatrix{
&& P \ar[ld]_-{p} \ar[rd]^-{q}  \\
& B_2 \ar[ld]_-{d_0} \ar[rd]^-{d_1} && B_2 \ar[ld]_-{d_1} \ar[rd]^-{d_2}  \\
B_1 && B_1  && B_1 \ ,}
$$
in which the diamond is a pullback, to obtain the graph
$$
\xymatrix{
 & P  \ar@<+1ex>[rr]|-{d_0 p} \ar@<-1ex>[rr]|-{d_2 q}  &&   B_{1} \ ,
}
$$
which underlies our category $B^{\mathbf{2}}$. There is a graph isomorphism
\begin{equation}\label{arrobj}
\CC(U,B^{\mathbf{2}}) \cong \CC(U,B)^{\mathbf{2}}
\end{equation}
natural in $U$. The right-hand side of \eqref{arrobj} has a category structure: the arrow category of $\CC(U,B)^{\mathbf{2}}$.
Since pullbacks exist in $\CC$, we can construct the objects required to complete $B^{\mathbf{2}}$
to a simplicial object in $\CC$ and, by Yoneda, we transport the morphisms across \eqref{arrobj} then internalize them to $\CC$. Using Corollary \ref{discretedensity} or \eqref{nerveB}, we extend the isomorphism \eqref{arrobj} to an isomorphism   
\begin{equation}\label{arrobj2}
(\mathrm{Cat}\CC)(A,B^{\mathbf{2}}) \cong (\mathrm{Cat}\CC)(U,B)^{\mathbf{2}} \ ,
\end{equation}
2-natural in $A$. This says that $B^{\mathbf{2}}$ is the cotensor of $\mathbf{2}$ and $B$ in the 2-category $\mathrm{Cat}\CC$.  

The 2-category $\mathrm{Cat}\CC$ has pullbacks formed pointwise with the simplicial objects. We can therefore compose the span
$$
\xymatrix{
B \ar@{<-}[r]^-{d_0}  & B^{\mathbf{2}} \ar[r]^-{d_1}  & B  }
$$
with itself; for example, we obtain
$$
\xymatrix{
&& B^{\mathbf{3}} \ar[ld]_-{p} \ar[rd]^-{q}  \\
& B^{\mathbf{2}} \ar[ld]_-{d_0} \ar[rd]^-{d_1} && B^{\mathbf{2}} \ar[ld]_-{d_1} \ar[rd]^-{d_2}  \\
B && B  && B }
$$
and then a functor $d_1:B^{\mathbf{3}} \lra B^{\mathbf{2}}$ lifting $d_1:B_2 \lra B_1$. This gives a category
$$
\xymatrix{
\mathrm{sq}B \ : \ B^{\mathbf{3}} \ar@<+2ex>[rr]|-{d_0} \ar@<+0ex>[rr]|-{d_1} \ar@<-2 ex>[rr]|-{d_2} &&   B^{\mathbf{2}}  \ar@<+2ex>[rr]|-{d_0} \ar@{<-}[rr]|-{i} \ar@<-2ex>[rr]|-{d_1} && B  }
$$
in $\mathrm{Cat}\CC$ such that there is a natural isomorphism
$$ \CC(U,\mathrm{sq}B)\cong \mathrm{sq}\CC(U,B)$$
where the right-hand side is explained at \eqref{sqdoubcat}. 
Using \eqref{nerveB}, we extend the isomorphism to
\begin{equation}
(\mathrm{Cat}\CC)(A,\mathrm{sq}B)\cong \mathrm{sq}(\mathrm{Cat}\CC)(U,B) \ .
\end{equation}

Looking at the fibration
$$\mathrm{cod} : (\mathrm{Cat}\CC)^{\mathbf{2}} \lra \mathrm{Cat}\CC$$
and the category $\mathrm{sq}B$ in $\mathrm{Cat}\CC$, we obtain a category
$$\mathrm{Spl}_{\mathrm{op}B} : = ((\mathrm{Cat}\CC)^{\mathbf{2}})^{\mathrm{sq}B} \ .$$
By Proposition \ref{FibasEM}, we have an equivalence  
\begin{equation}
\mathrm{Spl}_{\mathrm{op}B} \simeq ((\mathrm{Cat}\CC)/B)^{T_{\mathrm{sq}B}} \ .
\end{equation}
The monad $T_{\mathrm{sq}B}$ is easily identified. For any functor $f:X \lra B$ in $\CC$, form the pullback as in the square
$$
\xymatrix{
f/B \ar[rr]^-{q} \ar[d]_-{d_0} && B^{\mathbf{2}} \ar[d]^-{d_0} \ar[rr]^-{d_1}  && B \\
X \ar[rr]_-{f} && B \ .}
$$
The endofunctor of $T_{\mathrm{sq}B}$ takes $f:X\lra B$ to $d_1 = d_1 q : f/B \lra B$. 
We see that $T_{\mathrm{sq}B}$  is actually a 2-monad on $(\mathrm{Cat}\CC)/B$ 
yielding a 2-category $\mathrm{Spl}_{\mathrm{op}B}$. 

The objects of $\mathrm{Spl}_{\mathrm{op}B}$ are called \textit{split opfibrations over} $B$ in $\CC$.
Such an object is a functor $p:E\lra B$ in $\CC$ equipped with a functor $d_1:p/B\lra E$ for which
there exists a category structure in $(\mathrm{Cat}\CC)^{\mathbf{2}}$ extending the following diagram.
\begin{equation}\label{Splopfib}
\xymatrix{
p/B \ar[d]_-{} \ar@<+2ex>[rr]|-{d_{0}} \ar@{<-}[rr]|-{i} \ar@<-2ex>[rr]|-{d_{1}} &&E \ar[d]^-{p} \\
B^{\mathbf{2}} \ar@<+2ex>[rr]|-{d_{0}} \ar@{<-}[rr]|-{i} \ar@<-2ex>[rr]|-{d_{1}} && B}
\end{equation}
The functor $i$ in the top line of \eqref{Splopfib} is defined by the requirement that the pasted composite
2-cell
\begin{equation}
\xymatrix{
E \ar@/_/[ddr]_1 \ar@/^/[drrr]^p
\ar@{-->}[dr]^{i} \\
& p/B \ar[d]_{d_0}^(0.5){\phantom{aaa}}="1" \ar[rr]^{d_1}
&& B \ar[d]^{1}_(0.5){\phantom{aaa}}="2" \ar@{=>}"1";"2"^-{\lambda} \\
& E \ar[rr]_p && B }
\end{equation}
should be the identity 2-cell of $p$. In fact, the action $d_1:p/B\lra E$ of the monad $T_{\mathrm{sq}B}$
is unique up to isomorphism since one can show that it is left adjoint $d_1\dashv i$ to $i$.

We already defined discrete opfibration in Section \ref{IFSC}. Replacing $C$ by $B$ in the pullback \eqref{discopfib},
we can deduce the pullback 
\begin{equation}\label{discopfibint}
\xymatrix{
E^{\mathbf{2}} \ar[rr]^-{d_0} \ar[d]_-{p^{\mathbf{2}}} && E \ar[d]^-{p} \\
B^{\mathbf{2}} \ar[rr]_-{d_0} && B}
\end{equation}
in $\mathrm{Cat}\CC$. Define $j:E^{\mathbf{2}} \lra p/B$ by the following equation.
\begin{equation}
\xymatrix{
E^{\mathbf{2}} \ar@/_/[ddr]_{d_0} \ar@/^/[drrr]^{pd_1}
\ar@{-->}[dr]^{j} \\
& p/B \ar[d]_{d_0}^(0.5){\phantom{aaa}}="1" \ar[rr]^{d_1}
&& B \ar[d]^{1}_(0.5){\phantom{aaa}}="2" \ar@{=>}"1";"2"^-{\lambda} \\
& E \ar[rr]_p && B } \quad
\xymatrix{
\\ =} \quad
 \xymatrix{\\ E^{\mathbf{2}}  \ar@/^1pc/[rr]^{d_0} 
\ar@/_1pc/[rr]_{d_1} \ar@{}[rr]|{\Downarrow \; \lambda}&& E  \ar[rr]^p && B}
\end{equation}
It follows that a functor $p:E\lra B$ in $\CC$ is a discrete opfibration if and only if $j$ is invertible. We
can see then that $p$ is indeed an opfibration by taking the $d_1$ of \eqref{Splopfib} to be the composite
\begin{equation}
p/B\stackrel{j^{-1}}{\lra} E^{\mathbf{2}} \stackrel{d_1}{\lra} E \ .
\end{equation}
We can also then see that every commutative triangle
\begin{equation}
\xymatrix{
E \ar[rd]_{p}\ar[rr]^{f}   && F \ar[ld]^{q} \\
& B  &
}
\end{equation}
in $\mathrm{Cat}\CC$, with $q$ a discrete opfibration and $p$ any split opfibration, is a morphism in $\mathrm{Spl}_{\mathrm{op}B}$. Write $\mathrm{Disc}_{\mathrm{op}B}$ for the full subcategory of $\mathrm{Spl}_{\mathrm{op}B}$ whose objects are the discrete 
opfibrations. Our observations show that the 2-functor
\begin{equation}
\mathrm{Disc}_{\mathrm{op}B}\stackrel{\mathrm{incl.}}{\lra} \mathrm{Spl}_{\mathrm{op}B} \stackrel{{\mathrm{und.}}}{\lra} (\mathrm{Cat}\CC)/B
\end{equation}
is fully faithful.

Let $p:E\lra C$ be a discrete opfibration in $\CC$. Recall the definition of $\bar{p}^U :\CC(U,C)\lra \CC/U$ in Section~\ref{IFSC} \eqref{pbp}. For each category $B$ in $\CC$, we have a morphism of diagrams
\begin{equation}
\xymatrix{
\ar@<+2ex>[rr]|-{\CC(d_0,1)} \ar@{<-}[rr]|-{\CC(i,1)} \ar@<-2ex>[rr]_(0){\phantom{.}}="2a"|-{\CC(d_2,1)} \CC(B_0,C) \ar@{}[ddrr]|-{ } && \CC(B_1,C) \ar@<+2ex>[rr]|-{\CC(d_0,1)} \ar[rr]|-{\CC(d_1,1)} \ar@<-2ex>[rr]_(0){\phantom{.}}="1a"_(1){\phantom{.}}="0a"|-{\CC(d_2,1)} \ar@{}[ddrr]|-{ } && \CC(B_2,C)  \\ \\
\ar@<+2ex>[rr]^(0){\phantom{.}}="2b"|-{d_{0}^{\star}} \ar@{<-}[rr]|-{i^{\star}} \ar@<-2ex>[rr]|-{d_{1}^{\star}} \CC/B_0 && \CC/B_1 \ar@<+2ex>[rr]^(0){\phantom{.}}="1b"^(1){\phantom{.}}="0b"|-{d_{0}^{\star}} \ar[rr]|-{d_{1}^{\star}} \ar@<-2ex>[rr]|-{d_{2}^{\star}} && \CC/B_2 
\ar"2a";"2b"_{\bar{p}^{B_0}} \ar"1a";"1b"_{\bar{p}^{B_1}} \ar"0a";"0b"_{\bar{p}^{B_2}}
}\end{equation}
in the `pseudo' sense. This induces a functor
\begin{equation}\label{catpbp}
\bar{p}^B :(\mathrm{Cat}\CC)(B,C) \lra \mathrm{Spl}_{\mathrm{op}B}
\end{equation}
on the descent categories. Clearly then:

\begin{proposition} If $p:E\lra C$ is an internal full subcategory of $\CC$ then, for all categories $B$ in $\CC$, the functors $\bar{p}^B$ of (\ref{catpbp}) are fully faithful.
\end{proposition}

\begin{corollary}\label{2toposofintcats} 
If $\CC$ is a finitely complete, cartesian closed category and $p:E\lra C$ is an internal full subcategory then the 2-category $\mathrm{Cat}\CC$ equipped with the usual duality involution interchanging $d_0$ and $d_1$, and the classifying discrete opfibration $p$, is a 2-topos in the sense of Weber \cite{Weber2007}.
\end{corollary} 

\section{Size and cocompleteness}\label{Sac}

In Section 5 of \cite{Weber2007}, given a 2-topos, it is shown how to construct a 
Yoneda structure in the sense of Street-Walters \cite{StW1978} which is `good' in Weber's sense (compare Theorem 7 of \cite{EC}). After Corollary~\ref{2toposofintcats}, we are in a position to do this. 
However, it is of particular interest when the Yoneda morphisms provide some kind
of cocompletion of their domain object. This section addresses that question by
looking at cocompleteness concepts for an internal full subcategory. 

If $p:E\lra C$ is an internal full subcategory of $\CC$, we say a morphism $q:F\lra U$ in $\CC$ is $C$\textit{-fibred} when there exist a morphism $f:U\lra C$ and a pullback square
\begin{equation}
\xymatrix{
F \ar[rr]^-{ } \ar[d]_-{q} && E \ar[d]^-{p} \\
U \ar[rr]_-{f} && C & .}
\end{equation}  
Clearly $C$-fibred morphisms are stable under pullback.   

We say the internal full subcategory $C$ \textit{has coproducts} when the composite of any two composable $C$-fibred morphisms is $C$-fibred. We say the internal full subcategory $C$ \textit{has a terminator} when every identity morphism $1_U:U\lra U$ in $\CC$ is $C$-fibred.  

An object $A$ of any 2-category $\CK$ is said to \textit{admit coequalizers} 
when, for all objects $X$, the category $\CK(X,A)$ admits coequalizers and 
these are preserved by all functors of the form 
$\CK(h,A):\CK(X,A)\lra \CK(Y,A)$ where $h:Y\lra X$. Let $\mathbf{Pp}$ be the free category on the `parallel pair' directed graph
$$
\xymatrix{
 & \centerdot  \ar@<+1ex>[rr]|-{} \ar@<-1ex>[rr]|-{}  &&   \centerdot \ .
}
$$
A functor from $\mathbf{Pp}$ to a category $\CA$ amounts to a pair 
of morphisms in $\CA$ with the same domain and the same codomain. 
Suppose the cotensor $A^{\mathbf{Pp}}$ of $\mathbf{Pp}$ with $A$ exists
in the 2-category $\CK$. There is a `diagonal morphism' 
$\delta : A \lra A^{\mathbf{Pp}}$ corresponding to the parallel pair 
$(1_A,1_A)$ of morphisms in the category $\CK(A,A)$. It is easy to see that
the object $A$ admits coequalizers if and only if the morphism $\delta$ 
has a left adjoint.

\begin{proposition}\label{lanintofsc}
Suppose $C$ is an internal full subcategory of the finitely complete 
category $\CC$. The following properties pertain to the 2-category $\CK = \mathrm{Cat}\CC$.
\begin{itemize}
\item[(i)] If $C$ has a terminator then each category $\CC(U,C)$ has a terminal object preserved by each functor $\CC(r,C): \CC(U,C)\lra \CC(V,C)$ for $r:V\lra U$. Indeed, there exists a right adjoint $t$ to the unique functor $C\lra 1$.
\item[(ii)] If $C$ has coproducts and $h:U\lra V$ in $\CC$ is $C$-fibred 
then every morphism 
$f:U\lra C$ has a pointwise left extension along $h$. 
\item[(iii)] If $C$ has coproducts and admits coequalizers, and $h:A\lra B$ is a 
functor for which $h_0:A_0\lra B_0$, $d_1:A_1\lra A_0$ and $d_1:B_1\lra B_0$ 
are $C$-fibred, then every functor $f:A\lra C$ has a left extension along $h$. 
\end{itemize} 
\end{proposition} 

\begin{proof}
\begin{itemize}
\item[(i)] By hypothesis, there is an object $u$ in $\CC(U,C)$ taken by the fully faithful $\bar{p}^U$ of \eqref{pbp} to the terminal object $1_U$ of $\CC/U$. It follows that $u$ is a terminal object of $\CC(U,C)$. By pseudonaturality of the $\bar{p}^U$, the functor $\CC(r,C)$ preserves terminal objects since pullback $r^{\star}:\CC/U\lra \CC/V$ along $r$ does. A right adjoint for $C\lra 1$ is any $t:1\lra C$ with $\bar{p}^1(t)=1_1$.      
\item[(ii)] Since $h^{\star}:\CC/V\lra \CC/U$ has a left adjoint $\Sigma_h$ 
defined by composition with $h$, and since each $\bar{p}^U(f)$ is $C$-fibred, $C$ 
having coproducts implies $\Sigma_h(\bar{p}^U(f))$ is in the image of $\bar{p}^V$. 
So $\CC(h,1):\CC(V,C)\lra \CC(U,C)$ has a left adjoint obtained by restricting
$\Sigma_h$ along the components of $\bar{p}$. Now next notice that each slice
of the form $h/b$ is actually a pullback 
$$
\xymatrix{
P \ar[rr]^-{s } \ar[d]_-{r} && X \ar[d]^-{b} \\
U \ar[rr]_-{h} && V & .}
$$
in $\mathrm{Cat}\CC$ since $V$ is in $\CC$. 
Form the cube   
$$
\xymatrix{
\CC(V,C) \ar[rrr]^-{\CC(h,1)} \ar[rd] \ar[ddd]_-{\CC(b,1)} & & & \CC(U,C) \ar[ld] \ar[ddd]^-{\CC(r,1)} \\
& \CC/V \ar[d]_-{b^{\star}} \ar[r]^-{h^{\star}} & \CC/U \ar[d]^-{r^{\star}} \\
& \mathrm{Disc}_{\mathrm{op}X} \ar[r]_-{s^{\star}} & \mathrm{Disc}_{\mathrm{op}P} \\
\CC(X,C) \ar[ru] \ar[rrr]_-{\CC(s,1)} & & & \CC(P,C) \ar[lu] } 
$$
in which all face squares commute up to isomorphism (indeed, the big front 
face commutes on the nose). The sloping inward edges are all fully faithful
since they are components of $\bar{p}$. We need to see that the big front face
commutes when the top and bottom edges are replaced by their left adjoints.
Since $U$ and $V$ are discrete objects of $\mathrm{Cat}\CC$, the morphism $h$
is a discrete opfibration. So the pullback $s$ of $h$ along $b$ is also a discrete opfibration. 
So the left adjoint of $s^{\star}$ is the functor $\Sigma_s$ defined by composition 
with $s$. Since a horizontally pasted composite of pullbacks is a pullback, the square
$$
\xymatrix{
\CC/U \ar[rr]^-{\Sigma_s} \ar[d]_-{r^{\star}} && \CC/V \ar[d]^-{b^{\star}} \\
\mathrm{Disc}_{\mathrm{op}P} \ar[rr]_-{\Sigma_h} && \mathrm{Disc}_{\mathrm{op}X}}
$$
commutes up to isomorphism. Evaluating this square at a $C$-fibred object $t:F\lra U$ 
of $\CC$ and using that $ht$ is $C$-fibred, we obtain an object of 
$\mathrm{Disc}_{\mathrm{op}X}$ in the replete image of $\bar{p}^X$, as required.   
  
\item[(iii)]  Consider commutative the square
$$
\xymatrix{
\CK(B,C) \ar[rr]^-{\CK(h,1)} \ar[d]_-{} && \CK(A,C) \ar[d]^-{} \\
\CC(B_0,C) \ar[rr]_-{\CC(h_0,1)} && \CC(A_0,C)  .}
$$
It follows from Example \ref{codescofdiscrete} and Proposition \ref{FibasEM} that the 
left and right sides of the square are monadic since $\CC(d_1,1_C)$ has a left adjoint
for both $d_1:A_1\lra A_0$ and $d_1:B_1\lra B_0$.   Since $\CK(B,C)$ has
coequalizers and $\CC(h_0,1)$ has a left adjoint, the adjoint triangle theorem 
of Dubuc \cite{Dubuc1968} implies that $\CK(h,1)$ has a left adjoint, as required.  \hfill{$\square$}
\end{itemize} 
\end{proof} 

Using ideas of Theorems 3 and 28 of \cite{EC}, I strongly suspect that the left extensions in (iii) of Proposition \ref{lanintofsc} are pointwise. At this point, a direct proof eludes me. 

\begin{proposition}\label{Eastslice}
Suppose $C$ is an internal full subcategory of the finitely complete 
category $\CC$. Suppose $C$ has a terminator with $t:1\lra C$ right adjoint to the unique $C\lra 1$. Then there exists a 2-cell
\begin{equation}
\xymatrix{
E \ar[d]_{}^(0.5){\phantom{aaa}}="1" \ar[rr]^{p}  && C \ar[d]^{1_C}_(0.5){\phantom{aaa}}="2" \ar@{=>}"1";"2"^-{\lambda}
\\
1 \ar[rr]_-{t} && C 
}
\end{equation}  
in $\mathrm{Cat}\CC$ exhibiting $E$ as the slice $t/C$. Moreover, the morphism $t:1\lra C$ of $\mathrm{Cat}\CC$ is dense \eqref{dense}. 
\end{proposition}

\begin{proof} First note that the fully faithful functor $\bar{p}^X$ takes $X\lra 1 \stackrel{t}\lra C$ to $1_X:X\lra X$. So $\bar{p}^X$ on morphisms determines a bijection between 2-cells 
$$  
\xymatrix{
X \ar[d]_{}^(0.5){\phantom{aaa}}="1" \ar[rr]^{f}  && C \ar[d]^{1_C}_(0.5){\phantom{aaa}}="2" \ar@{=>}"1";"2"^-{\theta}
\\
1 \ar[rr]_-{t} && C 
}
$$
and morphisms $X \lra \bar{p}^X(f)$ over $X$. Since $ \bar{p}^X(f)$ is a pullback of $p$ along $f$, these are in bijection with morphisms $g:X\lra E$ such that $pg=f$. Taking $\lambda$ to correspond to $g = 1_E$ and a bit more work with 2-cells, we obtain the slice property required.

For the last sentence of the Proposition, we need to see that the diagram
$$  
\xymatrix{
t/b \ar[d]_{!}^(0.5){\phantom{aaa}}="1" \ar[rr]^{}  && X \ar[d]^{b}_(0.5){\phantom{aaa}}="2" \ar@{=>}"1";"2"^-{\lambda}
\\
1 \ar[rr]_-{t} && C 
}
$$
exhibits $b$ as a left extension of $t!$ along the top horizontal morphism. By the universal property of the slice $t/f$, 2-cells 
$$  
\xymatrix{
t/b \ar[d]_{!}^(0.5){\phantom{aaa}}="1" \ar[rr]^{}  && X \ar[d]^{f}_(0.5){\phantom{aaa}}="2" \ar@{=>}"1";"2"^-{\theta}
\\
1 \ar[rr]_-{t} && C 
}
$$
are in bijection with morphisms $t/b\lra t/f$ over $X$. Since the fully faithful $\bar{p}^X$ is defined by slicing $t/-$, these morphisms $t/b\lra t/f$ over $X$ are in bijection with 2-cells $\phi:b\Lra f:X\lra C$, as required. \hfill{$\square$}
\end{proof}

Proposition~\ref{Eastslice} suggests a simplification of the description of an internal full subcategory with terminator. Recall Definitions \ref{ff} and \ref{dense}.

\begin{proposition} Suppose $t:1\lra C$ is a fully faithful dense morphism in $\mathrm{Cat}\CC$ where $\CC$ is a finitely complete category. Then $C$ together with the discrete opfibration $d_1:t/C\lra C$ is an internal full subcategory of $C$ with terminator.
\end{proposition} 

\begin{proof} Since $t$ is dense, the square
$$  
\xymatrix{
t/f \ar[d]_{!}^(0.5){\phantom{aaa}}="1" \ar[rr]^{d_1}  && U \ar[d]^{f}_(0.5){\phantom{aaa}}="2" \ar@{=>}"1";"2"^-{\lambda}
\\
1 \ar[rr]_-{t} && C 
}
$$
exhibits $f$ as a left extension of $t!$ along $d_1$. It follows that morphisms $\theta:f\lra g$ in the category $\CC(U,C)$ are in bijection with 2-cells $\phi:t!\lra gd_1$. But such $\phi$ are in bijection with morphisms $t/f \lra t/g$ over $U$ by the universal property of the slice $t/g$. So the functor $\CC(U,C)\lra \CC/U$, taking $f$ to $d_1:t/f \lra U$, is fully faithful. That functor is pullback along $d_1:t/C\lra C$. It follows that $C$ is an internal full subcategory as asserted.  

To say $t:1\lra C$ is fully faithful is to say the square
$$  
\xymatrix{
1 \ar[d]_{}^(0.5){\phantom{aaa}}="1" \ar[rr]^{}  && 1 \ar[d]^{t}_(0.5){\phantom{aaa}}="2" \ar@{=}"1";"2"^-{}
\\
1 \ar[rr]_-{t} && C 
}
$$
has the slice property for $t/t$. So the square
$$  
\xymatrix{
U \ar[d]_{1_U}^(0.5){\phantom{aaa}}="1" \ar[rr]^{}  && 1 \ar[d]^{t}_(0.5){\phantom{aaa}}="2" \ar@{=}"1";"2"^-{}
\\
U \ar[rr]_-{t!} && C 
}
$$
has the slice property for $t!/t$. It follows that $C$ has a terminator.  \hfill{$\square$}
\end{proof}


\section{Internal full subcategories of $\mathrm{Cat}$}\label{IFSCoCat}

An algorithm for finding internal full subcategories of a locally presentable category $\CC$ was provided in \cite{CoIC}. This was applied to $\CC = \mathrm{Cat}$ in 
Section 8.10 of that paper. However, here we shall give an
internal full subcategory of $\mathrm{Cat}$ without going through the discovery process.

For a (small) category $A$, write $\hat{A}$ for the category $[A^{\mathrm{op}},\mathrm{Set}]$ of contravariant set-valued functors on $A$. It is the small-colimit completion of $A$ in the sense that 
restriction along the Yoneda embedding $y_A:A\lra \hat{A}$ yields an equivalence of categories
\begin{equation}\label{cocompletion}
\mathrm{Cocts}(\hat{A},X)\simeq [A,X] \ ,
\end{equation}
where the left side is the full subcategory of $[\hat{A},X]$ consisting of the small-colimit-preserving functors into the small cocomplete category $X$.

Let $\mathrm{mod_0}$ denote the category whose objects are small categories $A$ and whose morphisms $m:A\lra B$ are colimit-preserving functors $m:\hat{A} \lra \hat{B}$. Each morphism
$m:A\lra B$ determines a functor 
$$\overline{m}:B^{\mathrm{op}}\times A\lra \mathrm{Set}$$
defined by
\begin{equation}
\overline{m}(b,a)=m(A(-,a))(b) \ .
\end{equation}
By \eqref{cocompletion}, $m$ is uniquely determined up to isomorphism by $\overline{m}$. For
each $\phi \in A$, we have a natural family
$$\kappa : \phi(a)\times \overline{m}(b,a) \lra m(\phi)(b)$$
defined by
\begin{equation}
\kappa(x,y)=m(\hat{x})_b(y)
\end{equation}
where $\hat{x}:A(-,a)\lra \phi$ is the unique natural transformation with
$$\hat{x}_a(1,a)=x \in \phi(a) \ . $$
In particular, for $m:A\lra B$ and $n:B \lra C$ in $\mathrm{mod_0}$, we have
$$\kappa :\overline{m}(b,a)\times \overline{n}(c,b) \lra \overline{nm}(c,a) \ ,$$
which is a universal extraordinary natural family in the variable $b$; in other words, $\kappa$ induces an isomorphism
\begin{equation}\label{modcompcoend}
\int^b \overline{m}(b,a)\times \overline{n}(c,b) \cong \overline{nm}(c,a) \ .
\end{equation}

Let $T:X\lra \mathrm{mod_0}$ be a functor and define a category $\mathrm{El(T)}$ as follows. The
objects are pairs $(x,a)$ where $x\in X$ and $a \in Tx$. Morphisms $(\xi ,\tau):(x,a)\lra (y,b)$ 
consist of $\xi:x\lra y$ in $X$ and $\tau \in \overline{T\xi}(b,a)$. Composition is defined by
\begin{equation}\label{Elcomp}
\xymatrix{
(x,a) \ar[rd]_{(\xi , \tau)}\ar[rr]^{(\zeta \xi , \upsilon \star \tau)}   && (z,c) \ar@{<-}[ld]^{(\zeta , \upsilon)} \\
& (y,b)  &
}
\end{equation}
where $ \upsilon \star \tau$ is the value of the function
$$\kappa : \overline{T\xi}(b,a)\times \overline{T\zeta}(c,b) \lra \overline{T(\zeta \xi)}(c,a) $$
at $(\tau,\upsilon)$. The identity morphism of $(x,a)$ is $(1_x,1_a)$ where we note that
$$\overline{T1_x}(a,a)=(Tx)(a,a) \ .$$
There is a projection functor 
\begin{equation}\label{Elproj}
p_T:\mathrm{El(T)}\lra X
\end{equation}
defined by $p_T(x,a)=x$ and $p_T(\xi,\tau)=\xi$.

Let $\mathrm{mod_1}$ denote the category whose objects are functors $f:A\lra B$ between small categories. A morphism is a square
\begin{equation}\label{modsquare}
\xymatrix{
A \ar[d]_{m}^(0.5){\phantom{aaa}}="1" \ar[rr]^{f}  && B \ar[d]^{n}_(0.5){\phantom{aaa}}="2" \ar@{=>}"1";"2"^-{\theta}
\\
C \ar[rr]_-{g} && D 
}
\end{equation}
where $f$ and $g$ are the domain and codomain functors, where $m$ and $n$ are morphisms 
of $\mathrm{mod_0}$, and where $\theta$ is a natural family of functions
$$\theta_{b,a} : \overline{m}(c,a) \lra  \overline{n}(gc,fa) \ .$$
There is a bijection between such $\theta$ and natural transformations $\theta$ as in the diagram
\begin{equation}\label{hatsquare}
\xymatrix{
\hat{A} \ar[d]_{m}^(0.5){\phantom{aaa}}="1" \ar@{<-}[rr]^{\hat{f}}  && \hat{B} \ar[d]^{n}_(0.5){\phantom{aaa}}="2" \ar@{=>}"1";"2"^-{\theta}
\\
\hat{C} \ar@{<-}[rr]_-{\hat{g}} && \hat{D} & .
}
\end{equation}
Composition is achieved by vertically pasting the squares of the form \eqref{hatsquare}.

We have functors
\begin{equation}\label{gphmod}
d_0, d_1 : \mathrm{mod_1}\lra \mathrm{mod_0}
\end{equation}
defined by 
$$d_0(f)=A, \ \ d_0(m,\theta ,n)=m,$$
$$d_1(f)=B, \ \ d_1(m,\theta ,n)=n,$$
referring to \eqref{modsquare}. This defines a graph
in $\mathrm{Cat}$. There is also a canonical structure of category in $\mathrm{Cat}$
having \eqref{gphmod} as underlying graph. Composition for this category $\mathrm{mod}$
in $\mathrm{Cat}$ is derived from horizontal pasting of squares of the form \eqref{hatsquare}.   

At present, we are regarding $\mathrm{Cat}$ as a category, not a 2-category. When we write 
$\mathrm{Cat}(X,\mathrm{mod})$ we mean in the sense of 
$\CC(U,C)$ as described at the beginning of Section \ref{Cic}.

Now we shall extend the construction of \eqref{Elproj} to a functor 
\begin{equation}\label{El}
\mathrm{El} : \mathrm{Cat}(X,\mathrm{mod})\lra \mathrm{Cat}/X \ .
\end{equation}
Let $T$ and $S$ be objects of $\mathrm{Cat}(X,\mathrm{mod})$; that is, they are functors from $X$ to $\mathrm{mod_0}$. A morphism $\Theta:T\lra S$ in $\mathrm{Cat}(X,\mathrm{mod})$ is 
a functor $\Theta : X \lra \mathrm{mod}$ with $d_0 \Theta = T$ and $d_1 \Theta = S$. Each
morphism $\xi : x\lra y$ in $X$ yields a diagram
\[
\xymatrix{
Tx \ar[d]_{T\xi}^(0.5){\phantom{aaa}}="1" \ar[rr]^{\Theta_x}  && Sx \ar[d]^{S\xi}_(0.5){\phantom{aaa}}="2" \ar@{=>}"1";"2"^-{\theta_{\xi}}
\\
Ty \ar[rr]_-{\Theta_y} && Sy 
}
\]
in which $\Theta_x$ and $\Theta-y$ are functors. Define a functor
\begin{equation}\label{Elonarr}
\xymatrix{
\mathrm{El}(T) \ar[rd]_{p_T}\ar[rr]^{\mathrm{El}(\Theta)=F}   && \mathrm{El}(S) \ar[ld]^{p_T} \\
& X  &
}
\end{equation}
over $X$ as follows:
$$F(x,a)=(x,\Theta_x(a)), \ \ F(\xi,\tau)=(\xi,\theta_{\xi}(\tau))$$
using $\theta_{\xi} : \overline{T\xi}(b,a)\lra \overline{S\xi}(\Theta_y(b),\Theta_x(a))$. By
looking at \eqref{Elcomp} and using the fact that $\Theta$ is a functor, we 
see that $F$ is a functor. Clearly \eqref{Elonarr} commutes. Composition 
in $\mathrm{Cat}(X,\mathrm{mod})$ involves horizontal pasting of squares
\eqref{hatsquare}, from which we see that we do have a functor \eqref{El}.

The following result is related to Gray's Yoneda-like lemma on page 290 of \cite{Gray1969}; also see page 210 of Kelly \cite{Kelly1974a}.

\begin{theorem} The functors $\mathrm{El}$ of \eqref{El} are fully faithful for all $X$. The 
family of these functors is pseudonatural in $X$.
\end{theorem}

\begin{proof} 
The proof of the first sentence is a mere re-tracing of the steps in the definition of $\mathrm{El}$ on morphisms. The second sentence follows from the observation that we have
a pullback square
$$\xymatrix{
\mathrm{El}(Tr) \ar[rr]^-{} \ar[d]_-{p_{Tr}} && \mathrm{El}(T) \ar[d]^-{p_T} \\
Y \ar[rr]_-{r} && X}
$$
for all functors $r:Y\lra X$. \hfill{$\square$}
\end{proof}

\begin{corollary} The family of functors \eqref{El} exhibits $\mathrm{mod}$ as an internal full subcategory of $\mathrm{Cat}$.
\end{corollary}

By the generalized Yoneda lemma of Section \ref{FgYL} (see Theorem \ref{gYl} and \eqref{discopfib}), the pseudonatural family \eqref{El} is determined by a discrete opfibration  
\begin{equation}\label{genericp}
p: \mathrm{obj}\lra \mathrm{mod}
\end{equation}
between categories in $\mathrm{Cat}$. We shall describe $\mathrm{obj}$ explicitly.

The category $\mathrm{obj_0}$ has objects pairs $(A,a)$ where $A$ is a small category
and $a$ is an object of $A$. A morphism $(m,\mu):(A,a)\lra (B,b)$ 
in $\mathrm{obj_0}$
consists of a morphism $m:A\lra B$ in $\mathrm{mod_0}$ together 
with $\mu \in \overline{m}(b,a)$. Composition is a special case of \eqref{Elcomp} and uses the $\kappa$ of \eqref{modcompcoend}. We have the projection functor
\begin{equation}\label{genericp0}
p_0: \mathrm{obj_0}\lra \mathrm{mod_0}
\end{equation}
taking $(m,\mu):(A,a)\lra (B,b)$ to $m:A\lra B$.

\begin{proposition} The functor $p_0$ of \eqref{genericp0} is 
powerful in the category $\mathrm{Cat}$. That is, the functor 
$$p_0^{\star} : \mathrm{Cat}/\mathrm{mod}_0 \lra \mathrm{Cat}/\mathrm{obj}_0$$
has a right adjoint.
\end{proposition}

\begin{proof}
We must show that the functor $p_0$ has the factorization lifting property of Giraud-Conduch\'e (see \cite{Giraud1964}, \cite{Cond1972} and \cite{StVer}). Take a composable pair
$$A\stackrel{m}\lra B\stackrel{n}\lra C$$
in $\mathrm{mod_0}$ and a lifting
$$(A,a)\stackrel{(nm,\lambda)}\lra (C,c)$$
to $\mathrm{obj_0}$ of the composite $nm$. By \eqref{modcompcoend},
there exists $b\in B$ and $(\mu,\nu) \in \overline{m}(b,a)\times \overline{n}(c,b)$ 
such that $\kappa (\mu,\nu)=\lambda$. This gives a factorization
$$(A,a)\stackrel{(m,\mu)}\lra (B,b)\stackrel{(n,\nu)}\lra (C,c)$$
of $(nm,\lambda)$ which, using \eqref{modcompcoend} again, 
determines a unique path component in the category of such 
liftings, as required.      \hfill{$\square$}
\end{proof}

Now we shall define the category $\mathrm{obj}_1$. The objects are pairs $(f,a)$ where $f:A\lra B$ is a functor and $a \in A$. A morphism
\begin{equation}\label{obj1morph}
(m,\mu,n,\theta):(f,a)\lra (g,c)
\end{equation}
consists of a morphism $(m,n,\theta):f\lra g$ as in \eqref{modsquare} and 
$\mu \in \overline{m}(c,a)$. Composition is such that we have the functor 
\begin{equation}\label{genericp1}
p_1: \mathrm{obj_1}\lra \mathrm{mod_1}
\end{equation}
taking \eqref{obj1morph} to $(m,n,\theta)$. We also have the functor
\begin{equation}\label{genericd0}
d_0: \mathrm{obj_1}\lra \mathrm{obj_0}
\end{equation}
taking \eqref{obj1morph} to $(m,\mu):(A,a)\lra (C,c)$, and the functor 
\begin{equation}\label{genericd1}
d_1: \mathrm{obj_1}\lra \mathrm{obj_0}
\end{equation}
taking \eqref{obj1morph} to $(n,\theta_{c,a}(\mu)):(B,fa)\lra (D,gc)$. By
the general theory in Section \ref{IFSC}, the functors $d_0$ and $d_1$
of \eqref{genericd0} and \eqref{genericd1} form the underlying graph
in $\mathrm{Cat}$ of a category $\mathrm{obj}$ in $\mathrm{Cat}$, and
the functors $p_0$ and $p_1$ of \eqref{genericp0} and \eqref{genericp1}
form a functor $$p:\mathrm{obj}\lra \mathrm{mod}$$ in $\mathrm{Cat}$. Indeed,
we also know that $p$ is a discrete opfibration, pullback along which gives 
the functor $\mathrm{El}$ of \eqref{El}. 

Categories in $\mathrm{Cat}$ are called \textit{double categories}. They were defined by 
Ehresmann in \cite{Ehres1963}. Another reference is \cite{KelSt1974}. We write $\mathrm{Dbl}$ for the 2-category $\mathrm{Cat}(\mathrm{Cat})$ of 
double categories. Actually, looking only at this 2-category structure on
$\mathrm{Dbl}$ loses quite a bit of the symmetry of the situation. Since
$\mathrm{Cat}$ is cartesian closed, so too is $\mathrm{Dbl}$ and that is
important from the viewpoint of the 2-topos structure. It also means that
$\mathrm{Dbl}$ is hom-enriched in itself. There are two underlying functors
$$\mathrm{Dbl} \lra \mathrm{Cat}$$
which induce two 2-functors
$$\mathrm{Dbl\text{-}Cat} \lra \mathrm{Cat\text{-}Cat}=\mathrm{2\text{-}Cat}$$
whose values at $\mathrm{Dbl}$ itself give two ways to regard $\mathrm{Dbl}$
as a 2-category.



\begin{center}
--------------------------------------------------------
\end{center}

\appendix
\begin{thebibliography}{000}

\bibitem{Ben1967} Jean B\'enabou, \textit{Introduction to bicategories}, Lecture Notes in Mathematics \textbf{47} (Springer-Verlag, 1967) 1--77.\label{Ben1967}

\bibitem{Ben1973Top} Jean B\'enabou, \textit{Probl\`emes dans les topos}, Univ. Catholique de Louvain, Inst. de Math. Pure et Appliqu\'ee, Rapport no. \textbf{34} (1973).\label{Ben1973Top} 

\bibitem{Ben1975Fib} Jean B\'enabou, \textit{Fibrations petites et localement petites}, Comptes Rendus de l'Acad\'emie des Sciences, Paris, \textbf{281} (1975) 897--900.\label{Ben1975Fib} 

\bibitem{BenRou1970} Jean B\'enabou and Jacques Roubaud, \textit{Monades et descente}, Comptes Rendus de l'Acad\'emie des Sciences, Paris, \textbf{270} (1970) 96--98.\label{BenRou1970} 

\bibitem{BCSW} Renato Betti, Aurelio Carboni, Ross Street and Robert Walters, \textit{Variation through enrichment}, Journal of Pure and Applied Algebra \textbf{29} (1983) 109--127.\label{BCSW}

\bibitem{Borceux} Francis Borceux, \textit{Handbook of Categorical Algebra 1, 2 and 3}, Encyclopedia of Mathematics and its Applications \textbf{50}, \textbf{51} and \textbf{52} (Cambridge University Press, 1994).\label{Borceux}

\bibitem{CKS} Aurelio Carboni, Stefano Kasangian and Ross Street, \textit{Bicategories of spans and relations}, Journal of Pure and Applied Algebra \textbf{33} (1984) 259--267.\label{CKS}

\bibitem{Cond1972} Fran\c cois Conduch\'e, Au sujet de l'existence d'adjoints \`a droite aux foncteurs image r\'eciproque dans la cat\'egorie des cat\'egories, Compte Rendue Acad. Sci. Paris S\'er. A-B \textbf{275} (1972) A891--A894.\label{Cond1972}

\bibitem{DayConv} Brian J. Day, \textit{On closed categories of functors}, Lecture Notes in Mathematics \textbf{137} (Springer-Verlag, 1970) 1--38.\label{DayConv}

\bibitem{Dubuc1968} Eduardo Dubuc, \textit{Adjoint triangles}, Lecture Notes in Mathematics \textbf{61} (Springer-Verlag 1968) 69--91.\label{Dubuc1968}

\bibitem{Ehres1963} Charles Ehresmann, \textit{Cat\'egories doubles et cat\'egories structur\'ees} Comptes Rendue Acad. Sci. Paris \textbf{256} (1963) 1198--1201.\label{Ehres1963} 

\bibitem{EilKel1966} Samuel Eilenberg and G. Max Kelly, \textit{Closed categories}, Proceedings of the Conference on Categorical Algebra (La Jolla, 1965), (Springer-Verlag,1966) 421--562.\label{EilKel1966}

\bibitem{Giraud1964} Jean Giraud, \textit{M\'ethode de la descente}, Bull. Soc. Math. France M\'em. \textbf{2} (1964).\label{Giraud1964}

\bibitem{GPS} Robert Gordon, A. John Power and Ross Street, \textit{Coherence for tricategories}, Memoirs of the American Mathematical Society \textbf{117} no. 558 (1995) vi+81 pp.\label{GPS}

\bibitem{Gray1969} John W. Gray, \textit{The categorical comprehension scheme}, Lecture Notes in Mathematics \textbf{99} (Springer-Verlag, Berlin and New York, 1969) 242--312.\label{Gray1969}

\bibitem{GrothDesc} Alexander Grothendieck, \textit{Technique de descente et th\'eor\`emes d'existence en g\'eom\'etrie alg\'ebrique. I. G\'en\'eraliti\'es. Descente par morphismes fid\`element plats} S\'eminaire Bourbaki \textbf{5} (Expos\'e 190; 1959) viii+150 pp.\label{GrothDesc}

\bibitem{GrothFond} Alexander Grothendieck, \textit{Rev\^etements \'etales et groupe fondamental} \textbf{Fasc. II}
(Expos\'es 6, 8 \`a 11; S\'eminaire de G\'eom\'etrie Alg\'ebrique, 1960/61. Troisi\`eme \'edition, corrig\'ee) Institut des Hautes \'Etudes Scientifiques, Paris 1963 i+163 pp.\label{GrothFond}

\bibitem{Johnstone1977} Peter T. Johnstone, \textit{Topos theory} London Mathematical Society Monographs \textbf{10} (Academic Press, London-New York, 1977) xxiii+367 pp.\label{Johnstone1977}

\bibitem{JohnstWraith} Peter T. Johnstone and Gavin C. Wraith, \textit{Algebraic theories in toposes}, Lecture Notes in Mathematics \textbf{661} (Springer, Berlin, 1978) 141--242.\label{JohnstWraith}

\bibitem{Kan1958} Daniel M. Kan, \textit{Adjoint functors}, Transactions of the American Mathematical Society \textbf{87} (1958) 294--329. \label{Kan1958}

\bibitem{Kelly1969} G. Max Kelly, \textit{Adjunction for enriched categories}, Lecture Notes in Mathematics \textbf{106} (Springer-Verlag, 1969) 166--177.\label{Kelly1969}

\bibitem{Kelly1974a} G. Max Kelly, \textit{On clubs and doctrines}, Lecture Notes in Mathematics \textbf{420} (Springer-Verlag, 1974) 181--256.\label{Kelly1974a}

\bibitem{Kelly1974b} G. Max Kelly, \textit{Doctrinal adjunction}, Lecture Notes in Mathematics \textbf{420} (Springer-Verlag, 1974) 257--280.\label{Kelly1974b}

\bibitem{KellyBook} G. Max Kelly, \textit{Basic concepts of enriched category theory}, London Mathematical Society Lecture Note Series \textbf{64} (Cambridge University Press, Cambridge, 1982). \label{KellyBook}

\bibitem{KelSt1974} G. Max Kelly and Ross Street, \textit{Review of the elements of 2-categories}, Lecture Notes in Mathematics \textbf{420} (Springer-Verlag, 1974) 75--103.\label{KelSt1974}

\bibitem{Lack2000} Stephen Lack, \textit{A coherent approach to pseudomonads}, Advances in Mathematics \textbf{152} (2000) 179--202.\label{Lack2000} 

\bibitem{FTMII} Stephen Lack and Ross Street, \textit{The formal theory of monads II}, Journal of Pure and Applied Algebra \textbf{175} (2002) 243--265.\label{FTMII} 

\bibitem{Law1969} F. William Lawvere, \textit{Ordinal sums and equational doctrines},  Lecture Notes in Mathematics \textbf{80} (Springer-Verlag, 1969) 141--155.\label{Law1969}

\bibitem{CWM} Saunders Mac Lane, \textit{Categories for the Working Mathematician}, Graduate Texts in Mathematics \textbf{5} (Springer-Verlag, 1971).\label{CWM} 

\bibitem{MacLPar} Saunders Mac Lane and Robert Par\'e, \textit{Coherence for bicategories and indexed categories}, Journal of Pure and Applied Algebra \textbf{37} (1985) 59--80.\label{MacLPar}

\bibitem{Marmolejo} Francisco Marmolejo, \textit{Distributive laws for pseudomonads}, Theory and Applications of Categories \textbf{5} (1999) 91--147.\label{Marmolejo}

\bibitem{Par\'e2011} Robert Par\'e, \textit{Yoneda theory for double categories}, Theory and Applications of Categories \textbf{25} (2011) 436--489.\label{Par\'e2011} 

\bibitem{FTM} Ross Street, \textit{The formal theory of monads}, Journal of Pure and Applied Algebra \textbf{2} (1972) 149--168.\label{FTM} 

\bibitem{FYl} Ross Street, \textit{Fibrations and Yoneda's lemma in a 2-category}, Lecture Notes in Mathematics \textbf{420} (Springer-Verlag, 1974) 104--133.\label{FYl}

\bibitem{EC} Ross Street, \textit{Elementary cosmoi I}, Lecture Notes in Mathematics \textbf{420} (Springer-Verlag, 1974) 134--180.\label{EC}

\bibitem{Licv2f} Ross Street, \textit{Limits indexed by category-valued 2-functors}, Journal of Pure and Applied Algebra \textbf{8} (1976) 149--181.\label{Licv2f}

\bibitem{FiB} Ross Street, \textit{Fibrations in bicategories}, Cahiers de topologie et g\'eom\'etrie diff\'erentielle \textbf{21} (1980) 111--160.\label{FiB}

\bibitem{CoIC} Ross Street, \textit{Cosmoi of internal categories}, Transactions of the American Mathematical Society \textbf{258} (1980) 271--318.\label{CoIC}

\bibitem{HCSCSE} Ross Street, \textit{Higher categories, strings, cubes and simplex equations}, Applied Categorical Structures \textbf{3} (1995) 29--77.\label{HCSCSE}

\bibitem{CatStr} Ross Street, \textit{Categorical structures}, Handbook of Algebra \textbf{1} (edited by M. Hazewinkel, Elsevier Science, 1996) 529--577.\label{CatStr}

\bibitem{StVer} Ross Street and Dominic Verity, \textit{The comprehensive factorization and torsors}, Theory and Applications of Categories \textbf{23} (2010) 42--75.\label{StVer}

\bibitem{StW1978} Ross Street and Robert F.C. Walters, \textit{Yoneda structures on 2-categories}, Journal of Algebra \textbf{50} (1978) 350--379.\label{StW1978} 

\bibitem{Weber2007} Mark Weber, \textit{Yoneda structures from 2-toposes}, Applied Categorical Structures \textbf{15} (2007) 259--323.\label{Weber2007}

\bibitem{Wraith1975} Gavin C. Wraith, \textit{Lectures on elementary topoi}, Lecture Notes in Mathematics, \textbf{445} (Springer, Berlin, 1975) 114--206.\label{Wraith1975}

\bibitem{Zoeb} Volker Z\"oberlein, \textit{Doctrines on 2-categories}, Mathematische Zeitschrift \textbf{148} (1976) 267--279. \label{Zoeb}



\end{thebibliography}

\end{document}